% concatenated from main.tex
\documentclass[12pt, reqno]{amsart}
\usepackage[foot]{amsaddr}

\usepackage[mathscr]{eucal}
\usepackage[all]{xy}
\usepackage{graphicx}
\usepackage{caption}
\usepackage{subcaption}
\captionsetup[figure]{name=Figure, labelfont=normalfont}
\captionsetup[subfigure]{labelfont=normalfont}
\usepackage{multicol, multirow}
\usepackage{amssymb,amsfonts,amsthm}
\usepackage{color}
\usepackage[centertags]{amsmath}
\usepackage{empheq}
\usepackage{geometry}
\usepackage{algorithm}
\usepackage{algpseudocode}
\usepackage{hyperref}
\usepackage{mathrsfs}
\usepackage{enumerate}
\usepackage{esint}
\usepackage{IEEEtrantools}
\usepackage{booktabs}
\usepackage{float}

\pagestyle{myheadings}
\textwidth=16cm
\textheight=23cm
\parindent=16pt
\oddsidemargin=-0.5cm
\evensidemargin=-0.5cm
\topmargin=-0.5cm

\theoremstyle{plain}
\newtheorem{thm}{Theorem}[section]
\newtheorem{prop}[thm]{Proposition}
\newtheorem{lem}[thm]{Lemma}
\newtheorem{conclusion}[thm]{Conclusion}

\theoremstyle{definition}
\newtheorem{definition}[thm]{Definition}
\newtheorem{remark}[thm]{Remark}

\numberwithin{equation}{section}
\numberwithin{figure}{section}
\numberwithin{algorithm}{section}

\title[]{Hamiltonian Interface Dynamics for Reduced-Order Optimization of Incompressible Mixing}

\author{Ziqian Li$^1$$^2$}
\address{$^1$Chair for Dynamics, Control, Machine Learning and Numerics, Department of Mathematics, Friedrich-Alexander-Universit\"{a}t Erlangen-N\"{u}rnberg, 91058 Erlangen, Germany}
\address{$^2$School of Mathematics, Jilin University, Changchun, Jilin 130012, China}


\author{Enrique Zuazua$^1$$^3$$^4$}
\address{$^3$Chair of Computational Mathematics, Deusto University, 48007 Bilbao, Basque Country, Spain}
\address{$^4$Departamento de Matem\'aticas, Universidad Aut\'onoma de Madrid, 28049 Madrid, Spain}

\email{ziqian.li@fau.de, enrique.zuazua@fau.de}

\begin{document}

\keywords{%
Fluid mixing, optimal control, optimal design, Hamiltonian dynamics, model reduction, interface length, symplectic integration.}


\begin{abstract}
We develop a reduced-order framework for optimizing mixing in two-dimen-sional incompressible flows. Instead of optimizing the full transport PDE, the method maximizes the length of advected material interfaces, leading to a finite-dimensional Hamiltonian control problem based on parametrized stream functions. We derive the continuous adjoint equations and reduced gradients, and discretize the forward and adjoint dynamics with the implicit midpoint rule. The resulting discrete adjoint is algebraically consistent with the derivative of the fully discrete objective, up to the tolerance of the nonlinear midpoint solves. The approach applies to bounded two-dimensional domains with smooth finite-dimensional stream-function parametrizations. Numerical experiments on cellular-flow and Doswell frontogenesis benchmarks show that the optimized time-dependent Hamiltonians generate near-exponential interface stretching and substantially faster decay of the $\dot{H}^{-1}$ mix-norm, in contrast with the polynomial behavior observed for stationary flows. When evaluated on a common reference transport solver, the interface-based controls produce faster $\dot{H}^{-1}$ decay than a Eulerian Sobolev-norm optimizer under a matched setup, while substantially reducing computational cost. We also identify a limitation of the reduced model: increasing the control basis may further improve the interface-length objective without yielding proportional gains in $\dot{H}^{-1}$ mixing, confirming that interface length is an effective but not fully faithful proxy for mixing in geometrically complex regimes.
\end{abstract}

\maketitle

\section{Introduction}

\subsection{Problem statement and motivation}\label{sec:problem}
Understanding and controlling mixing by incompressible transport is a central issue in fluid mechanics, scientific computing, and applications ranging from microfluidics to geophysical transport. In this paper, we consider a passive scalar $\theta$ transported by an incompressible velocity field $\mathbf v$ in a bounded planar domain $\Omega\subset\mathbb R^2$. The governing model is the transport equation
\begin{equation}\label{eq:transport}
\partial_t\theta+\mathbf v\cdot\nabla\theta=0
\qquad\text{in }(0,T)\times\Omega,
\end{equation}
together with the incompressibility and no-penetration conditions
\begin{equation}\label{eq:incomp_bc}
\nabla\cdot\mathbf v=0
\qquad\text{in }(0,T)\times\Omega,
\qquad
\mathbf v\cdot\mathbf n_\Omega=0
\qquad\text{on }(0,T)\times\Gamma,
\end{equation}
where $\Gamma:=\partial\Omega$ and $\mathbf n_\Omega$ is the outward unit normal. The \emph{optimal mixing problem} is to choose the stirring field $\mathbf v$ so that the transported scalar becomes mixed as efficiently as possible.

A large part of the modern literature quantifies mixing by negative Sobolev norms, often called \emph{mix-norms}. The homogeneous $\dot{H}^{-1}(\Omega)$ norm
\begin{equation}\label{eq:mixnorm}
\|\theta\|_{\dot{H}^{-1}(\Omega)}
:=\sup\Bigl\{
\int_\Omega\theta\,\phi\,dx
\;:\;
\int_\Omega\phi\,dx=0,\;
\int_\Omega|\nabla\phi|^2\,dx\le 1
\Bigr\}
\end{equation}
is a natural measure of mixedness, as introduced by Mathew, Mezi\'c, and Petzold \cite{mathew2005multiscale} and widely employed in the analysis of mixing rates \cite{crippa2017cellular}. A smaller value of $\dot{H}^{-1}(\Omega)$ norm corresponds to better mixing, reflecting the transfer of scalar content from large to small spatial scales.

However, negative Sobolev norms are not the only meaningful way to quantify mixing. In many two-phase transport settings, geometric quantities such as interface deformation, geometric mixing scales, or entropy-type functionals also provide natural descriptions of the mixing process; see, for example, \cite{vikhansky2002enhancement, alberti2016exponential, thiffeault2012using}.

The main purpose of this paper is to develop such a geometric viewpoint in a Hamiltonian framework. We take the final length of the material interface as the objective functional to be optimized, rather than the final $\dot{H}^{-1}(\Omega)$ norm of the scalar. This choice is motivated by theoretical, computational, and modelling considerations, as we now explain.

In two space dimensions, an incompressible velocity field admits a stream-function representation $\mathbf v=\nabla^\perp \psi$, under which the characteristic dynamics becomes Hamiltonian. Denoting by $J=\bigl(\begin{smallmatrix}0&-1\\1&\phantom{-}0\end{smallmatrix}\bigr)$ the symplectic matrix, the flow map satisfies $\dot\Phi_\psi=J\nabla_x \psi(t,\Phi_\psi)$. An initial material interface $\gamma_0$ is propagated by $\Phi_\psi$, and its position $X(t,s)=\Phi_\psi(t;\gamma_0(s))$ and tangent vector $a(t,s)=\partial_s X(t,s)$ satisfy the coupled ODE system
\[
\partial_t X = J\nabla_x \psi(t,X),
\qquad
\partial_t a = J\nabla_x^2 \psi(t,X)\,a.
\]
The final length of the interface $\mathcal{L}(\psi;T)=\int_0^1|a(T,s)|\,ds$ is therefore expressed entirely in terms of the Hamiltonian flow and its differential, replacing the Eulerian full-field mixing objective by a Lagrangian geometric one, and the method yields a substantial computational reduction, quantified in Section~\ref{sec:comparison}.

The use of $\mathcal L(\psi;T)$ as a surrogate for mixing has a transparent geometric rationale. For a two-phase configuration whose transition layer is concentrated near the material interface $\Gamma_t$, the area balance between the two phases suggests the heuristic scaling
\begin{equation}\label{eq:geom_scale}
\ell(t)\,\mathcal L(\psi;t)\sim|\Omega|,
\end{equation}
where $\ell(t)$ denotes the typical filament width; growth of the interface length therefore naturally corresponds to decrease of a geometric mixing scale. This argument is heuristic and does not establish equivalence with negative Sobolev mix-norms; we use it solely to motivate the reduced objective, whose predictive value is assessed a posteriori by evaluating the optimized controls on the full transport equation~\eqref{eq:transport}.

The reduction underlying our framework is Hamiltonian/Lagrangian (replacing the full Eulerian transport PDE by ODEs on a small set of advected markers) rather than a multigrid hierarchy in the sense of Glowinski's two-grid philosophy \cite{Glowinski1992}; the spirit, however, is similar: optimization is carried out on a reduced model that captures the essential geometric mechanism, and the resulting controls are ultimately evaluated on the original transport equation (see Section~\ref{sec:comparison}). This viewpoint is also in line with the role of interfaces as primary geometric objects in level-set methods~\cite{OsherSethian1988} and active contour models~\cite{Caselles1997}.


\subsection{Main results and contributions}\label{sec:main_results}

Our contributions are twofold: methodological and computational. On the methodological side, we introduce a reduced optimal-control framework for incompressible mixing based on the geometry of an advected material interface. On the computational side, the reduction yields lower optimization cost and improved $\dot{H}^{-1}$ mix-norm decay (in the tested benchmarks) when the resulting controls are evaluated on the original transport equation.

More precisely, our main results are as follows.

\begin{itemize}
\item \textbf{A geometric reduced model for mixing based on Hamiltonian interface dynamics.}
We show that, in two dimensions, the advection of a material interface by an incompressible flow is naturally governed by a Hamiltonian system: the interface position $X$ and its tangent vector $a=\partial_s X$ satisfy a closed system of ODEs, and the final length of the interface is determined entirely by the Hamiltonian flow and its linearization (Proposition~\ref{prop:interface_ode}). 

\item \textbf{An adjoint-based optimal-control formulation on the reduced model.}
We formulate a regularized control problem that maximizes the final interface length $\mathcal L(\psi;T)$ under an $L^2$ penalty on the control amplitude (equivalent, via the Gram matrix of the basis, to an $L^2(0,T;L^2(\Omega))$ penalty on the stream function). Within a finite-dimensional smooth stream-function ansatz, we prove the existence of an optimal Hamiltonian (Theorem~\ref{prop:H3}), derive the continuous adjoint system, and obtain an explicit gradient formula in terms of adjoint variables transported along marker trajectories. This is done avoiding full-domain PDE solves during the optimization loop.

\item \textbf{A fully discrete optimization framework.}
We discretize the forward and adjoint systems by the implicit midpoint rule, which preserves the Hamiltonian structure and is self-adjoint. As a consequence, the discrete adjoint obtained by transposing the linearized forward map coincides with the midpoint scheme run backward in time, so that the resulting discrete gradient is algebraically consistent with the derivative of the discrete cost at the scheme level, up to the tolerance of the nonlinear solves used to implement the midpoint step. The optimization is then carried out with a Polak--Ribi\`ere conjugate-gradient method and Armijo line search (Algorithm~\ref{alg:optimize}).

\item \textbf{Validation on the full transport PDE.}
The reduced model tracks $O(N_p)$ marker points instead of $O(N_{\mathrm{cell}}^2)$ Eulerian degrees of freedom per time step; this Hamiltonian/Lagrangian reduction is the direct source of the per-iteration cost reduction reported in Section~\ref{sec:comparison}. On cellular-flow and Doswell benchmarks under matched conditions, the resulting controls yield faster $\dot{H}^{-1}$ decay than the Eulerian Sobolev-norm optimizer of \cite{hu2026structure} (Table~\ref{tab:optimizer_comparison}), with interface-length and mix-norm growth/decay well fitted by exponential rates over the simulated horizons.

\item \textbf{Interpretation.}
A plausible heuristic explanation for the observed advantage is that the reduced formulation focuses on large-scale interface stretching and is less sensitive to the fine-grid numerical effects that affect a PDE-level optimizer (Section~\ref{sec:comparison} reports the supporting evidence).

\item \textbf{Limitations of the reduced model.}
Interface length is a geometric proxy, not equivalent to the $\dot{H}^{-1}$ objective: the enlarged-basis experiment of Section~\ref{sec:N4} produces a $16\times$ longer interface yet only a moderately smaller $\dot{H}^{-1}$ rate, so the surrogate's predictive power must ultimately be assessed on the original transport equation.

\end{itemize}

\medskip
\noindent\textbf{Novelty and significance.}
To our knowledge, this is the first framework that combines a Hamiltonian marker/tangent reduced model of an advected interface, an adjoint-based optimization of interface length within a finite-dimensional stream-function ansatz, a symplectic discrete adjoint algebraically consistent with the discrete cost gradient, and a systematic matched-conditions validation of the resulting controls against an Eulerian Sobolev-norm optimizer. The method is not proposed as a universal surrogate for mix-norm optimization: Section~\ref{sec:N4} identifies regimes in which interface stretching ceases to be a faithful proxy for $\dot{H}^{-1}$ mixing.


\subsection{Related work}

Our work draws on three lines of research: the optimal control of fluid mixing, the geometric treatment of moving interfaces, and symplectic time integration for Hamiltonian systems.

\textbf{Optimal control of fluid mixing.}
Control of transport and mixing via pure advection has been extensively studied through open-loop design under physical constraints \cite{d1999control, vikhansky2002enhancement} and through finite-dimensional velocity ansatzes \cite{mathew2007optimal}. Our recent work \cite{hu2026structure} develops a structure-preserving Eulerian optimizer for the Sobolev-norm objective; the present work complements \cite{hu2026structure} by replacing the Eulerian objective with a Lagrangian interface-length formulation.

\textbf{Symplectic integration.}
The characteristic flow of a two-dimensional incompressible velocity field is Hamiltonian, making structure-preserving time integration \cite{HLW06} essential for long-time accuracy. The implicit midpoint rule we employ is symmetric and self-adjoint, ensuring consistency between the discrete forward and adjoint maps.

\subsection{Organization of the paper}
The remainder of the paper is organized as follows. Section~\ref{sec:Hamiltonian} introduces the Hamiltonian interface dynamics and derives the evolution equations for the interface and its tangent vector. Section~\ref{sec:ocp} formulates the length-based optimal control problem, derives the continuous adjoint system, and presents the reduced gradient. Section~\ref{sec:numerics} describes the numerical framework, including the symplectic ODE solver, the discrete-adjoint gradient computation, and the optimization algorithm. Numerical experiments are presented in Section~\ref{sec:numerics_exp}, including stationary benchmarks (Section~\ref{sec:stationary_flows}), optimized flows (Sections~\ref{sec:optimal_flows} and~\ref{sec:optimal_Doswell}), a comparison with the Eulerian optimizer (Section~\ref{sec:comparison}), and an investigation of control-space dimensionality (Section~\ref{sec:N4}). Conclusions and perspectives are given in Section~\ref{sec:conclusions}.

\section{Hamiltonian interface dynamics}\label{sec:Hamiltonian}

In this section, we exploit the stream function formulation of two-dimensional incompressible flows to derive the Hamiltonian structure of the characteristic dynamics. We then show that the advected interface and its tangent vector satisfy a closed ODE system (Proposition~\ref{prop:interface_ode}), whose solution determines the length of the transported material curve. Since the stream-function formulation is intrinsically two-dimensional, we work throughout in two space dimensions.

Let $\Omega \subset \mathbb{R}^2$ be a bounded simply connected domain with $C^2$ boundary $\Gamma$, and let $T>0$. Let
\[
\psi:[0,T]\times\Omega\to\mathbb{R}
\]
be a time-dependent stream function satisfying the homogeneous boundary condition
\[
\psi(t,\cdot)\big|_{\Gamma}=0
\qquad\text{for all } t\in[0,T].
\]
The velocity field is then defined by
\begin{equation}\label{eq:velocity}
\mathbf{v}(t,x):=\nabla^\perp \psi(t,x)
:=\bigl(-\partial_{x_2}\psi(t,x),\,\partial_{x_1}\psi(t,x)\bigr).
\end{equation}
By construction, $\mathbf v$ is divergence-free in $\Omega$ and tangent to $\Gamma$, so \eqref{eq:incomp_bc} is automatically satisfied.

For each $x\in\Omega$, let $\Phi_\psi(t;x)$ denote the flow map generated by $\mathbf v$:
\begin{equation*}
\frac{d}{dt}\Phi_\psi(t;x)=\mathbf{v}\bigl(t,\Phi_\psi(t;x)\bigr),
\qquad
\Phi_\psi(0;x)=x.
\end{equation*}
Using \eqref{eq:velocity}, this takes the Hamiltonian form
\begin{equation}\label{eq:hamiltonian_flow}
\frac{d}{dt}\Phi_\psi(t;x)=J\nabla_x \psi\bigl(t,\Phi_\psi(t;x)\bigr),
\qquad
\Phi_\psi(0;x)=x,
\end{equation}
where
\[
J:=
\begin{pmatrix}
0 & -1\\
1 & 0
\end{pmatrix}
\]
is the canonical symplectic matrix. Thus, the characteristic dynamics admits a natural Hamiltonian structure.

Let $\gamma_0\in W^{2,\infty}([0,1];\Omega)$ (explained in \ref{rem:gamma0_reg}) be a parameterization of an initial material interface. The advected interface is
\[
X(t,s):=\Phi_\psi\bigl(t;\gamma_0(s)\bigr),
\qquad
(t,s)\in[0,T]\times[0,1].
\]
Differentiating with respect to the interface parameter $s$, we define the tangent vector
\[
a(t,s):=\partial_s X(t,s).
\]

\begin{prop}[Interface dynamics]\label{prop:interface_ode}
The pair $(X,a)$ satisfies the coupled ODE system
\begin{subequations}\label{eq:state}
\begin{align}
\partial_t X(t,s) &= J\nabla_x \psi\bigl(t,X(t,s)\bigr),\label{eq:state_X}\\
\partial_t a(t,s) &= J\nabla_x^2 \psi\bigl(t,X(t,s)\bigr)\,a(t,s),\label{eq:state_a}\\
X(0,s) &= \gamma_0(s),\quad a(0,s) = \gamma_0'(s).
\end{align}
\end{subequations}
The length of the advected interface at time $t$ is
\begin{equation}\label{eq:length}
\mathcal{L}(\psi;t):=\int_0^1 |a(t,s)|\,ds.
\end{equation}
\end{prop}

\begin{proof}
Equation \eqref{eq:state_X} is the Hamiltonian ODE \eqref{eq:hamiltonian_flow} evaluated at $x=\gamma_0(s)$. Differentiating \eqref{eq:state_X} with respect to $s$ yields \eqref{eq:state_a}, since $a=\partial_s X$ and
\[
\partial_t a=\partial_s\partial_t X=\partial_s\bigl(J\nabla_x \psi(t,X)\bigr)=J\nabla_x^2 \psi(t,X)\,\partial_s X=J\nabla_x^2 \psi(t,X)\,a.
\]
The length formula \eqref{eq:length} follows from the standard definition applied to the parameterized curve $s\mapsto X(t,s)$.
\end{proof}

Hence, the geometry of the transported interface is completely determined by the Hamiltonian flow and its differential along trajectories. The functional $\mathcal{L}(\psi;T)$ is the geometric quantity used below as a reduced surrogate objective for mixing.


\section{Optimal control problem}\label{sec:ocp}

Building on the Hamiltonian interface dynamics established in Section~\ref{sec:Hamiltonian}, in particular the closed system \eqref{eq:state} for the marker pair $(X,a)$ that serves as the ODE constraint of the optimal control problem, we now formulate the optimization problem of maximizing the terminal interface length. To avoid unnecessary technical difficulties, we introduce a regularized cost functional, define the admissible control set, derive the continuous adjoint system, and obtain the reduced gradient formula within the finite-dimensional ansatz.

\subsection{Length-maximization problem}

Fix $\varepsilon>0$ and define the regularized Euclidean norm
\begin{equation}\label{eq:ell}
\ell_\varepsilon(z):=\sqrt{|z|^2+\varepsilon^2},
\qquad
\nabla \ell_\varepsilon(z)=\frac{z}{\sqrt{|z|^2+\varepsilon^2}},
\end{equation}
which is a smooth approximation to $|z|$, differentiable at $z=0$.

We restrict the stream function to a finite-dimensional subspace. Fix $N\in\mathbb N$ and a family of linearly independent basis functions $h_1,\dots,h_N$ satisfying the standing assumption
\begin{equation}\label{eq:basis_assumption}
h_k\in C^\infty(\overline\Omega)\cap H_0^1(\Omega),\qquad k=1,\dots,N,
\end{equation}
and set
\begin{equation}\label{eq:smooth_space}
\mathcal H_N:=\mathrm{span}\{h_1,\dots,h_N\}\subset C^\infty(\overline\Omega)\cap H_0^1(\Omega).
\end{equation}
Since each $h_k\in C^\infty(\overline\Omega)$, the basis constants
\begin{equation}\label{eq:basis_constants}
B^{(j)}:=\max_{1\le k\le N}\|h_k\|_{W^{j,\infty}(\Omega)},\qquad j=0,1,2,3,
\end{equation}
are finite.

Fix a radius $R_u>0$ and consider the associated set of admissible stream function profiles
\begin{equation}\label{eq:Uad}
\mathcal U_{\mathrm{ad}}^{R_u}
:=
\Bigl\{
\psi(t,x)=\sum_{k=1}^N u_k(t)\,h_k(x)\;;\;u\in L^2(0,T;\mathbb R^N),\ \|u\|_{L^2(0,T;\mathbb R^N)}\le R_u
\Bigr\}.
\end{equation}
The dimension of the admissible set $N$ being fixed, for every $j\in\{0,1,2,3\}$,
\begin{equation}\label{eq:pt_bound}
\|\psi(t,\cdot)\|_{W^{j,\infty}(\Omega)}\le\sqrt N\,B^{(j)}\,|u(t)|\qquad\text{for a.e. }t\in(0,T),
\end{equation}
where $|u(t)|$ denotes the Euclidean norm of $u(t)\in\mathbb R^N$. This yields, in particular, bounds for $\psi\in L^2(0,T;W^{3,\infty}(\Omega))$ and $\mathbf v:=\nabla^\perp\psi\in L^2(0,T;W^{2,\infty}(\Omega;\mathbb R^2))$ in terms of the $L^2(0,T;\mathbb R^N)$-norm of $u(t).$


By the linear independence of $\{h_k\}_{k=1}^N$, each $\psi\in\mathcal U_{\mathrm{ad}}^{R_u}$ is determined by a unique coefficient vector $u\in L^2(0,T;\mathbb R^N)$ with $\|u\|_{L^2(0,T;\mathbb R^N)}\le R_u$; we accordingly identify $\psi$ with $u$. The symmetric positive-definite Gram matrix
\begin{equation}\label{eq:gram}
M\in\mathbb R^{N\times N},\qquad M_{k\ell}:=(h_k,h_\ell)_{L^2(\Omega)},\qquad k,\ell=1,\dots,N,
\end{equation}
records the cross-correlations of the basis; a direct expansion yields, for every $u\in\mathbb R^N$ and $\psi=\sum_{k=1}^N u_k h_k$,
\begin{equation}\label{eq:gram_identity}
\|\psi\|_{L^2(\Omega)}^2=u^\top M u.
\end{equation}

For $\lambda>0$, we define the cost functional
\begin{equation}\label{eq:cost}
\mathcal{J}_{\lambda,\varepsilon}(\psi)
:=
-\int_0^1 \ell_\varepsilon\bigl(a(T,s)\bigr)\,ds
+\frac{\lambda}{2}\int_0^T \|\psi(t,\cdot)\|_{L^2(\Omega)}^2\,dt,
\qquad\psi\in\mathcal U_{\mathrm{ad}}^{R_u},
\end{equation}
where $a(T,s)$ and $\psi(t,\cdot)$ satisfies the state system  \eqref{eq:state}.
By \eqref{eq:gram_identity}, the penalty admits the equivalent coefficient representation $$\frac{\lambda}{2}\int_0^T u(t)^\top M u(t)\,dt.$$
The optimal control or design problem we consider is
\begin{equation}\label{eq:ocp}
\min_{\psi\in\mathcal{U}_{\mathrm{ad}}^{R_u}}\mathcal{J}_{\lambda,\varepsilon}(\psi).
\end{equation}
The negative sign in the first term reflects that we \emph{maximize} the final interface length while penalizing the control effort.

The interface $\Gamma_t$ optimized in \eqref{eq:ocp} is not arbitrary: it coincides with the material image of $\Gamma_0$ under the Hamiltonian flow, which in turn coincides with the zero level set of the scalar field used in the full PDE validation of Section~\ref{sec:comparison}. We record this identification as a separate proposition because it ties the Lagrangian object that is optimized to the Eulerian object on which the resulting controls are subsequently evaluated.

\begin{prop}[Transport of the reference interface]\label{prop:zero_contour}
Let $\theta_0\in C^1(\overline\Omega)$ be such that $\Gamma_0:=\{x\in\overline\Omega:\theta_0(x)=0\}$ is a regular level set, and let $\theta\in C^1([0,T]\times\overline\Omega)$ solve the transport equation \eqref{eq:transport} for the velocity field $\mathbf v=\nabla^\perp\psi$ associated with some $\psi\in\mathcal U_{\mathrm{ad}}^{R_u}$. Then, for every $t\in[0,T]$,
\begin{equation*}
\Gamma_t:=\{x\in\overline\Omega:\theta(t,x)=0\}=\Phi_\psi(t;\Gamma_0),
\end{equation*}
where $\Phi_\psi(t;\cdot)$ is the Hamiltonian flow map of $\psi$. In particular, the material interface optimized in \eqref{eq:ocp} is the zero contour of the transported scalar that is used in the full-PDE validation of Section~\ref{sec:comparison}.
\end{prop}

\begin{proof}
By the method of characteristics, $\theta(t,x)=\theta_0(\Phi_\psi(t;x)^{-1})$ for $x\in\overline\Omega$, where $\Phi_\psi(t;\cdot)^{-1}$ exists by the smoothness and incompressibility of $\mathbf v$ and is the inverse of $\Phi_\psi(t;\cdot)$. The condition $\theta(t,x)=0$ is therefore equivalent to $\Phi_\psi(t;x)^{-1}\in\Gamma_0$, that is, $x\in\Phi_\psi(t;\Gamma_0)$.
\end{proof}

\begin{remark}[Interface length as a critical geometric mix-norm]\label{rem:critical_mixnorm}
The use of the interface length \eqref{eq:cost} as a reduced objective admits a precise distributional interpretation as a geometric analogue of a negative Sobolev mix-norm, which complements the area-balance heuristic \eqref{eq:geom_scale}. Let $\Gamma\subset\Omega$ be a smooth material curve and denote by $\mu_\Gamma:=\mathcal H^1\!\lfloor_\Gamma$ the one-dimensional Hausdorff measure restricted to $\Gamma$. If $E\subset\Omega$ is one of the two phases bounded by $\Gamma=\partial E$, then the distributional gradient of the characteristic function $\chi_E$ is the vector-valued Radon measure $D\chi_E=-n_\Gamma\,\mathcal H^1\!\lfloor_\Gamma$, and consequently
\[
|D\chi_E|(\Omega)=\mathcal H^1(\Gamma)=\mathcal L(\Gamma);
\] the interface length is exactly the total variation of the singular gradient of a two-phase scalar.

This observation can be related to negative Sobolev norms. The line measure $\mu_\Gamma$ belongs to $H^{-s}(\Omega)$ for every $s>1/2$, with $s=1/2$ the critical exponent for codimension-one measures in two dimensions \cite[Ch.~12]{mattila1995geometry}. If $\rho_\delta$ is a standard mollifier and $\mu_\Gamma^\delta:=\rho_\delta\ast\mu_\Gamma$, then at the critical level a formal Plancherel computation gives the leading-order asymptotics: for a flat segment of length $L$ along the $x_1$-axis, $\widehat{\mu_\Gamma}(\xi_1,\xi_2)=2\sin(\xi_1L/2)/\xi_1$ is constant in $\xi_2$, and the angular average $$\int_{S^1}|\widehat{\mu_\Gamma}(r\hat\xi)|^2\,d\hat\xi\sim 2\pi L/r$$ for $r=|\xi|\to\infty$ produces, after integration against $|\xi|^{-1}|\widehat{\rho_\delta}|^2$ in polar coordinates, $\int_0^{1/\delta}(2\pi L/r)\,dr=2\pi L\log(1/\delta)+O(1)$; the same scaling holds for general smooth curves by a partition of unity along arclength, with curvature and nonlocal interactions contributing only to the $O_\Gamma(1)$ remainder. We thus obtain
\[
\|\mu_\Gamma^\delta\|_{\dot{H}^{-1/2}(\Omega)}^2=C\,\mathcal L(\Gamma)\,\log(1/\delta)+O_\Gamma(1)\qquad\text{as }\delta\to 0,
\]
with a constant $C>0$ depending only on the Fourier normalization and on the choice of mollifier. Hence the renormalized critical Sobolev size of the regularized interface measure recovers the interface length:
\[
\mathcal L(\Gamma)\sim\frac{1}{C\log(1/\delta)}\,\|\mu_\Gamma^\delta\|_{\dot{H}^{-1/2}(\Omega)}^2\qquad\text{as }\delta\to 0.
\]
In this sense, maximizing $\mathcal L(\psi;T)$ may be viewed as increasing an inverse geometric mixing scale, or equivalently the renormalized critical Sobolev size of the interface measure carried by the optimized material interface. We emphasize that this interpretation does \emph{not} imply equivalence with the Eulerian $\dot{H}^{-1}$ mix-norm \eqref{eq:mixnorm} of the transported scalar; rather, it provides a distributional rationale for using interface length as a Lagrangian geometric surrogate for mixing in two-phase configurations, complementing the heuristic of \eqref{eq:geom_scale}.
\end{remark}

We are now in the position to state the main well-posedness and existence result for \eqref{eq:ocp}.



\begin{thm}[Existence of a minimizer]\label{prop:H3}
Let $\mathcal H_N$, $B^{(j)}$, $\mathcal U_{\mathrm{ad}}^{R_u}$, and $\mathcal J_{\lambda,\varepsilon}$ be given by \eqref{eq:smooth_space}, \eqref{eq:basis_constants}, \eqref{eq:Uad}, and \eqref{eq:cost}, respectively. Assume that the basis $h_k$ satisfies \eqref{eq:basis_assumption} and that $\gamma_0\in W^{2,\infty}([0,1];\overline\Omega)$. Then the following properties hold.
\begin{enumerate}
\item[\textup{(i)}]\emph{(State regularity and well-posedness.)} For every $u\in L^2(0,T;\mathbb R^N)$ with associated stream function $\psi=\sum_{k=1}^N u_k h_k$,
\begin{equation}\label{eq:v_reg}
\psi\in L^2(0,T;W^{3,\infty}(\overline\Omega)),\qquad
\mathbf v=\nabla^\perp\psi\in L^2(0,T;W^{2,\infty}(\Omega;\mathbb R^2));
\end{equation}
the flow $\Phi_\psi$ maps $\overline\Omega$ into itself, and the state system  \eqref{eq:state} admits a unique solution $(X,a)\in C([0,T]\times[0,1];\overline\Omega\times\mathbb R^2)$. Consequently, the cost functional $\mathcal{J}_{\lambda,\varepsilon}$ is well defined and depends continuously on $u$, i.e., on  $\psi$.
\item[\textup{(ii)}]\emph{(Existence of a minimizer.)} The minimization problem \eqref{eq:ocp} admits at least one minimizer $\psi^*\in\mathcal U_{\mathrm{ad}}^{R_u}$.
\end{enumerate}
\end{thm}


\begin{proof}
The proof combines classical direct-method ingredients (Carath\'eodory, Gr\"onwall, Arzel\`a--Ascoli, lower semicontinuity) adapted to the finite-dimensional Ha-miltonian ansatz and the bounded set $\mathcal U_{\mathrm{ad}}^{R_u}$.

\noindent\textit{Proof of \textup{(i)}.}
The proof is organized in the following four steps.

\vskip 5pt
\noindent\textit{Step 1. Regularity of the stream function and velocity field.}

By \eqref{eq:pt_bound} with $j=3$, we readily get \eqref{eq:v_reg}.

\vskip 5pt
\noindent\textit{Step 2. Flow invariance of $\overline\Omega$.}

Since $\psi(t,\cdot)\in H_0^1(\Omega)\cap C^\infty(\overline\Omega)$ by \eqref{eq:basis_assumption} and \eqref{eq:smooth_space}, for every $t\in[0,T]$,
\begin{equation*}
\mathbf v(t,\cdot)\cdot\mathbf n_\Omega=\nabla^\perp\psi(t,\cdot)\cdot\mathbf n_\Omega=-\partial_\tau\psi(t,\cdot)=0\qquad\text{on }\Gamma,
\end{equation*}
where $\partial_\tau$ denotes the tangential derivative along $\Gamma$. Let $d:\mathbb R^2\to\mathbb R$ denote the signed distance to $\Gamma$, negative inside $\Omega$. Since $\Gamma$ is of class $C^2$, $d\in C^2$ in a tubular neighbourhood of $\Gamma$ with $\nabla d=\mathbf n_\Omega$ on $\Gamma$. For any absolutely continuous trajectory $X(\cdot)$ of \eqref{eq:state_X} with $X(t_0)\in\Gamma$,
\begin{equation}\label{eq:invariance_id}
\frac{d}{dt}d(X(t))\big|_{t=t_0}=\nabla d(X(t_0))\cdot\mathbf v(t_0,X(t_0))=0.
\end{equation}
By \eqref{eq:invariance_id} and the Lipschitz regularity of $\mathbf v$ from \eqref{eq:v_reg}, Nagumo's viability theorem \cite[Ch.~V, Thm.~4.1]{hartman2002ode} implies that $\overline\Omega$ is positively invariant under the flow $\Phi_\psi$.

\vskip 5pt
\noindent\textit{Step 3. Existence and uniqueness of $(X,a)$.}

Define $F(t,X,a;u):=\bigl(J\nabla_x\psi(t,X),\,J\nabla_x^2\psi(t,X)\,a\bigr)$. By the regularity considerations above, $F$ is measurable in $t$, of class $C^\infty$ in $(X,a)$, and locally Lipschitz in $(X,a)$ with constant
\begin{equation}\label{eq:F_lip}
L_F(t,|a|):=\sqrt N\,B^{(2)}\,|u(t)|+\sqrt N\,B^{(3)}\,|u(t)|\,(1+|a|)\in L^2(0,T)\subset L^1(0,T).
\end{equation}
Fix $s\in[0,1]$; since $\gamma_0\in W^{2,\infty}([0,1];\overline\Omega)$, the initial datum $(\gamma_0(s),\gamma_0'(s))$ is well defined. Carath\'eodory's existence theorem \cite[Ch.~II, Thm.~1.1]{hartman2002ode} produces a local solution, and \eqref{eq:F_lip} combined with Gr\"onwall's inequality yields uniqueness \cite[Ch.~II, Thm.~2.1]{hartman2002ode}.

It remains to exclude finite-time blow-up. Flow invariance (Step~2) confines $X(\cdot,s)$ to the compact set $\overline\Omega$. For $a$, the tangent equation \eqref{eq:state_a} and \eqref{eq:pt_bound} with $j=2$ give
\begin{equation*}
\frac{d}{dt}|a(t,s)|^2\le 2\,\|\nabla_x^2\psi(t,\cdot)\|_{L^\infty(\Omega)}\,|a(t,s)|^2\le 2\sqrt N\,B^{(2)}\,|u(t)|\,|a(t,s)|^2,
\end{equation*}
where Gr\"onwall's inequality yields the $u$-dependent bound
\begin{equation}\label{eq:Ma_u}
|a(t,s)|\le M_a(u):=\|\gamma_0'\|_{L^\infty([0,1])}\,\exp\!\Bigl(\sqrt N\,B^{(2)}\int_0^T|u(\tau)|\,d\tau\Bigr),
\end{equation}
which is finite for every $u\in L^2(0,T;\mathbb R^N)\subset L^1(0,T;\mathbb R^N)$. Hence the solution $(X,a)$ extends uniquely to all of $[0,T]$. Continuous dependence on $s$ follows from Gr\"onwall applied to two solutions sharing $u$, whose initial data differ by $(\|\gamma_0'\|_{L^\infty}+\|\gamma_0''\|_{L^\infty})|s_1-s_2|$; continuity in $t$ is automatic from the integral form of the ODE. Jointly, $(X,a)\in C([0,T]\times[0,1];\overline\Omega\times\mathbb R^2)$.

\vskip 5pt
\noindent\textit{Step 4. Length well-posedness.}

The regularized norm $\ell_\varepsilon$ is globally $1$-Lipschitz on $\mathbb R^2$ because $$|\nabla\ell_\varepsilon(z)|=|z|/\sqrt{|z|^2+\varepsilon^2}\le 1.$$ By Step~3, $a\in C([0,T]\times[0,1];\mathbb R^2)$ is uniformly continuous on the compact rectangle $[0,T]\times[0,1]$. Hence, for every $t,t'\in[0,T]$,
\begin{equation*}
|\mathcal L(\psi;t)-\mathcal L(\psi;t')|\le\int_0^1|\ell_\varepsilon(a(t,s))-\ell_\varepsilon(a(t',s))|\,ds\le\int_0^1|a(t,s)-a(t',s)|\,ds\xrightarrow{t'\to t}0,
\end{equation*}
so $\mathcal L(\psi;\cdot)\in C([0,T])$. This proves the last assertion of (i).

\vskip 5pt
\noindent\textit{Proof of \textup{(ii)}.}
We apply the direct method of the calculus of variations on $\mathcal U_{\mathrm{ad}}^{R_u}$, which by \eqref{eq:Uad} is itself the $L^2(0,T;\mathbb R^N)$-ball of radius $R_u$; the proof is organized in the following four steps.

\vskip 5pt
\noindent\textit{Step 1. A uniform \emph{a priori} bound on the tangent.}

For every $u\in L^2(0,T;\mathbb R^N)$ with $\|u\|_{L^2}\le R_u$, the Cauchy--Schwarz inequality yields $\int_0^T|u(\tau)|\,d\tau\le\sqrt T\,R_u$, so that the $u$-dependent bound \eqref{eq:Ma_u} of Step~3 of (i) specializes to
\begin{equation}\label{eq:Ma}
|a(t,s)|\le M_a:=\|\gamma_0'\|_{L^\infty([0,1])}\,\exp\,\bigl(\sqrt N\,B^{(2)}\,\sqrt T\,R_u\bigr)
\end{equation}
uniformly on $[0,T]\times[0,1]$ for every $\psi\in\mathcal U_{\mathrm{ad}}^{R_u}$. In particular,
\begin{equation*}
\mathcal J_{\lambda,\varepsilon}(\psi)\ge-\int_0^1\sqrt{|a(T,s)|^2+\varepsilon^2}\,ds\ge-\sqrt{M_a^2+\varepsilon^2},
\end{equation*}
so $\mathcal J_{\lambda,\varepsilon}$ is bounded below on $\mathcal U_{\mathrm{ad}}^{R_u}$.

\vskip 5pt
\noindent\textit{Step 2. Weak compactness of a minimizing sequence.}

Let $(\psi_m)\subset\mathcal U_{\mathrm{ad}}^{R_u}$ be a minimizing sequence with coefficients $u_m\in L^2(0,T;\mathbb R^N)$ satisfying $\|u_m\|_{L^2}\le R_u$. The closed ball of radius $R_u$ in Hilbert space $L^2(0,T;\mathbb R^N)$ is convex, strongly closed, and bounded, hence weakly compact; therefore, up to extraction, $u_m\rightharpoonup u^*$ weakly with $\|u^*\|_{L^2}\le R_u$. Setting $\psi^*(t,x):=\sum_{k=1}^N u^*_k(t)\,h_k(x)$, we have $\psi^*\in\mathcal U_{\mathrm{ad}}^{R_u}$.

\vskip 5pt
\noindent\textit{Step 3. Uniform convergence of the state.}

Let $(X_m,a_m)$ and $(X^*,a^*)$ be the states associated with $\psi_m$ and $\psi^*$, respectively. By Step~1, $|a_m(t,s)|\le M_a$ uniformly in $m,t,s$. We establish joint $(t,s)$-equicontinuity of $(X_m,a_m)$.

\emph{Equicontinuity in $t$.} Integrating \eqref{eq:state_X} and invoking the pointwise bound \eqref{eq:pt_bound} with $j=1$,
\begin{align*}
|X_m(t,s)-X_m(t',s)|
\le\int_{t'}^t\sqrt N\,B^{(1)}\,|u_m(\tau)|\,d\tau
\le\sqrt N\,B^{(1)}\,R_u\,\sqrt{|t-t'|}
\end{align*}
by Cauchy--Schwarz. Analogously, by \eqref{eq:state_a}, $|a_m|\le M_a$, and the pointwise bound \eqref{eq:pt_bound} with $j=2$,
\begin{align*}
|a_m(t,s)-a_m(t',s)|\le\sqrt N\,B^{(2)}\,M_a\,R_u\,\sqrt{|t-t'|}.
\end{align*}

\emph{Equicontinuity in $s$.} Since $\partial_s X_m(t,s)=a_m(t,s)$, we have $$|X_m(t,s_1)-X_m(t,s_2)| \le M_a|s_1-s_2|.$$ Differentiating the tangent ODE $\partial_t a_m=J\nabla_x^2\psi_m(t,X_m)\,a_m$ in $s$ yields
\begin{align*}
\partial_t(\partial_s a_m)
=J\nabla_x^2\psi_m(t,X_m)\,\partial_s a_m
+J\nabla_x^3\psi_m(t,X_m)[\,a_m,a_m\,],
\end{align*}
with initial datum $\partial_s a_m(0,s)=\gamma_0''(s)\in L^\infty([0,1])$. By the pointwise bound \eqref{eq:pt_bound} with $j=3$, the forcing is bounded by $\sqrt N\,B^{(3)}\,|u_m(t)|\,M_a^2$, while the coefficient of $\partial_s a_m$ has $L^\infty_x$-norm at most $\sqrt N\,B^{(2)}\,|u_m(t)|$. Gr\"onwall's inequality then produces a uniform bound
\begin{equation}\label{eq:Map}
|\partial_s a_m(t,s)|\le M_a':=\bigl(\|\gamma_0''\|_{L^\infty}+\sqrt N\,B^{(3)}\,M_a^2\,\sqrt T\,R_u\bigr)\,\exp\!\bigl(\sqrt N\,B^{(2)}\,\sqrt T\,R_u\bigr),
\end{equation}
and hence $|a_m(t,s_1)-a_m(t,s_2)|\le M_a'|s_1-s_2|$.

By Arzel\`a--Ascoli, a subsequence (not relabelled) satisfies $X_m\to\widetilde X$ and $a_m\to\widetilde a$ uniformly on $[0,T]\times[0,1]$. The smoothness of $h_k$ together with uniform convergence $X_m\to\widetilde X$ implies $J\nabla h_k(X_m)\to J\nabla h_k(\widetilde X)$ strongly in $L^2((0,T)\times(0,1);\mathbb R^2)$, while $u_{k,m}\rightharpoonup u^*_k$ weakly in $L^2(0,T)$. The product of a strongly $L^2$-convergent sequence and a weakly $L^2$-convergent sequence converges weakly in $L^1$, which suffices to pass the integral form of \eqref{eq:state_X}--\eqref{eq:state_a} to the limit; the quadratic term $J\nabla_x^2\psi_m(t,X_m)\,a_m$ is treated analogously. Uniqueness from (i) identifies $(\widetilde X,\widetilde a)=(X^*,a^*)$, so that $a_m(T,\cdot)\to a^*(T,\cdot)$ uniformly on $[0,1]$.

\vskip 5pt
\noindent\textit{Step 4. Lower semicontinuity and passage to the limit.}

By Step~3 and continuity of $\ell_\varepsilon$,
\begin{align*}
\int_0^1\ell_\varepsilon(a_m(T,s))\,ds\longrightarrow\int_0^1\ell_\varepsilon(a^*(T,s))\,ds.
\end{align*}
The penalty is convex and continuous on $L^2(0,T;\mathbb R^N)$: indeed, $u\mapsto\int_0^T u(t)^\top M u(t)\,dt$ is a continuous quadratic form, positive by the positive-definiteness of $M$ in \eqref{eq:gram}, hence weakly lower semicontinuous. Combining,
\begin{align*}
\mathcal J_{\lambda,\varepsilon}(\psi^*)\le\liminf_m\mathcal J_{\lambda,\varepsilon}(\psi_m)=\inf_{\mathcal U_{\mathrm{ad}}^{R_u}}\mathcal J_{\lambda,\varepsilon},
\end{align*}
so $\psi^*$ is a minimizer of \eqref{eq:ocp}.

\vskip 5pt
In summary, the proof of Theorem~\ref{prop:H3} is completed.
\end{proof}

Next we give three remarks on Theorem~\ref{prop:H3}: the first explains the regularity required of the initial interface~$\gamma_0$; the second explains why the ball constraint $\|u\|_{L^2}\le R_u$ built into $\mathcal U_{\mathrm{ad}}^{R_u}$ in \eqref{eq:Uad} cannot be dropped without losing existence; and the third justifies the $L^2(\Omega)$-penalty on~$\psi$.

\begin{remark}[Regularity of initial interface]\label{rem:gamma0_reg}
The hypothesis $\gamma_0\in W^{2,\infty}([0,1];\overline\Omega)$ is sharp for the argument of Theorem~\ref{prop:H3} and is precisely what matches the $W^{j,\infty}$-scale of the pointwise bound \eqref{eq:pt_bound}. Two derivatives enter the proof at distinct places:
\begin{enumerate}
\item \emph{Length measurement and a priori bounds on $a$.} The interface length \eqref{eq:length} and the bound \eqref{eq:Ma_u} both feature $\gamma_0'$ through the initial datum $a(0,s)=\gamma_0'(s)$; the hypothesis $\gamma_0'\in L^\infty([0,1])$ ensures that $\mathcal L(\psi;0)<\infty$ and that the Gr\"onwall-type estimate of Step~3 of~(i) produces the finite constant 
	\begin{align*}
		M_a(u)=\|\gamma_0'\|_{L^\infty}\exp(\sqrt N\,B^{(2)}\int_0^T|u|).
	\end{align*}
	Without $\gamma_0'\in L^\infty$, the objective $-\mathcal L(\psi;T)$ itself would fail to be defined on the admissible class.
\item \emph{Compactness of minimizing sequences.} The proof of existence in~(ii) rests on the uniform bound \eqref{eq:Map} on $\partial_s a_m$, whose initial datum is $\partial_s a_m(0,s)=\gamma_0''(s)$. The hypothesis $\gamma_0''\in L^\infty([0,1])$ makes this Gr\"onwall bound finite and hence furnishes the $s$-equicontinuity of $(X_m,a_m)$ needed to apply Arzel\`a--Ascoli in Step~2 of~(ii).
\end{enumerate}
Weaker regularity $\gamma_0\in W^{2,p}$ for some finite $p$ would preserve existence at the cost of propagating $L^p$-in-$s$ bookkeeping through the proof; the $W^{2,\infty}$-choice is the minimal one under which the pointwise bounds \eqref{eq:pt_bound} are inherited by $(X,a,\partial_s a)$ uniformly in $s$, keeping the analysis within a single Sobolev scale.
\end{remark}

\begin{remark}[On the role of control constraints]\label{rem:control_constraints}
We formulate the optimization problem under the constraint $\psi\in\mathcal U_{\mathrm{ad}}^{R_u}$ since, in the absence of such restrictions, the coercivity of the functional $\mathcal J_{\lambda,\varepsilon}$ is far from obvious. Indeed, the cost combines two competing effects: the maximization of the interface length and the penalization of the control effort. Because the Hamiltonian vector field advecting the interface depends linearly on the stream function, and hence on the control $u$, the Gr\"onwall-type estimate \eqref{eq:Ma_u} anticipates a possible exponential growth of the (negative) interface-length term in the norm of the control, which can dominate the positive quadratic control penalty. The dimension $N$ is fixed throughout this paper and plays no role in this competition; the delicacy concerns the dependence of the functional on the control parameters themselves, the infinite-dimensional design problem $N\to\infty$ lying outside the scope considered here. As a consequence, without suitable constraints, one cannot expect bounded minimizing sequences, and the existence of minimizers becomes unclear.

To make this competition precise at the level of the interface length itself, note that $\ell_\varepsilon(a)\le|a|+\varepsilon$ gives, after integration of \eqref{eq:Ma_u} over $s\in[0,1]$ and applying the Cauchy--Schwarz inequality $\int_0^T|u(\tau)|\,d\tau\le\sqrt T\,\|u\|_{L^2(0,T;\mathbb R^N)}$,
\begin{equation}\label{eq:length_u_bound}
\mathcal L(\psi;t)\le\|\gamma_0'\|_{L^\infty([0,1])}\,\exp\,\bigl(\sqrt{NT}\,B^{(2)}\,\|u\|_{L^2(0,T;\mathbb R^N)}\bigr)+\varepsilon,
\end{equation}
which grows exponentially in $\|u\|_{L^2}$, whereas the penalty in \eqref{eq:cost} grows only quadratically in $\|u\|_{L^2}$ (the Gram matrix $M$ of \eqref{eq:gram} has eigenvalues bounded above and below by positive constants). Consequently
\[
\mathcal J_{\lambda,\varepsilon}(\psi)\ge-\mathcal L(\psi;T)+\frac{\lambda}{2}\lambda_{\min}(M)\,\|u\|_{L^2}^2
\]
cannot be made to tend to $+\infty$ as $\|u\|_{L^2}\to\infty$: the functional $\mathcal J_{\lambda,\varepsilon}$ is \emph{not} globally coercive on the unbounded set $L^2(0,T;\mathbb R^N)$. Existence of a minimizer without any size constraint on $u$ is therefore an open problem, and the bounded-ball formulation \eqref{eq:Uad} provides a well-posed surrogate.

In Section~\ref{sec:numerics_exp}, with $R_u=10$, the terminal ratios $\|u^*\|_{L^2}/R_u$ across the five experiments are $\{0.16,\,0.15,\,1.00,\,0.70,\,0.15\}$: the ball is inactive in four cases, while the Doswell constant-initial-guess case saturates the constraint and exercises the projection step of Algorithm~\ref{alg:optimize}.
\end{remark}

\begin{remark}[Choice of the $L^2$ penalty]\label{rem:penalty_choice}
On $\mathcal H_N$, all norms are equivalent and the pointwise bound \eqref{eq:pt_bound} ensures that an $L^2$-in-time control on $u$ also controls $\psi$ in any required Sobolev norm. The $L^2(\Omega)$ penalty on $\psi$ is therefore the canonical choice: by \eqref{eq:gram_identity} it is equivalent to the coefficient-space penalty $\frac{\lambda}{2}\int_0^T u(t)^\top M u(t)\,dt$, and positive-definiteness of $M$ guarantees coercivity in $|u|$ on any bounded ball.
\end{remark}

\subsection{Adjoint system and reduced gradient}\label{sec:adjoint}

Having established well-posedness and existence in Theorem~\ref{prop:H3}, we now derive the first-order optimality system for \eqref{eq:ocp}. To dualize the state equations \eqref{eq:state}, we introduce a pair of adjoint variables $(p,q)$ and first verify that they are uniquely determined by a backward linear ODE.

\begin{lem}[Adjoint system and well-posedness]\label{lem:adjoint_wp}
Under the hypotheses of Theorem~\ref{prop:H3}, the backward linear system
\begin{subequations}\label{eq:adjoint}
\begin{align}
-\partial_t p(t,s)
&=J\nabla_x^2\psi(t,X)^\top p(t,s)+\bigl(J\nabla_x^3\psi(t,X)[\,\cdot\,,a(t,s)]\bigr)^{\!\top}\!q(t,s),\label{eq:adjoint_p}\\
-\partial_t q(t,s)
&=J\nabla_x^2\psi(t,X)^\top q(t,s),\label{eq:adjoint_q}\\
p(T,s)&=0,\qquad
q(T,s)=-\nabla\ell_\varepsilon(a(T,s))=-\frac{a(T,s)}{\sqrt{|a(T,s)|^2+\varepsilon^2}},\label{eq:adjoint_tc}
\end{align}
\end{subequations}
admits a unique solution $(p,q)\in C([0,T]\times[0,1];\mathbb R^2\times\mathbb R^2)$.
\end{lem}

\begin{proof}
For a.e.\ $t\in(0,T)$ and every $s\in[0,1]$, the pointwise bound \eqref{eq:pt_bound} (with $j=2,3$) together with $|a(t,s)|\le M_a$ from \eqref{eq:Ma} gives
\begin{align*}
\bigl|J\nabla_x^2\psi(t,X(t,s))\bigr|
&\le\|\nabla_x^2\psi(t,\cdot)\|_{L^\infty(\Omega)}\le\sqrt N\,B^{(2)}|u(t)|,\\
\bigl|J\nabla_x^3\psi(t,X(t,s))[\,\cdot\,,a(t,s)]\bigr|
&\le\|\nabla_x^3\psi(t,\cdot)\|_{L^\infty(\Omega)}\,|a(t,s)|
\le\sqrt N\,B^{(3)}\,M_a\,|u(t)|,
\end{align*}
so each coefficient matrix in \eqref{eq:adjoint_p}--\eqref{eq:adjoint_q} is bounded, uniformly in $s$, by a function of $t$ that lies in $L^2(0,T)\subset L^1(0,T)$ (since $u\in L^2(0,T;\mathbb R^N)$). The terminal datum \eqref{eq:adjoint_tc} is bounded, since $|\nabla\ell_\varepsilon|\le 1$. Equation~\eqref{eq:adjoint_q} is therefore a linear Carath\'eodory ODE for $q$ in $t$, parameterized by $s$, to which \cite[Ch.~II, Thms.~1.1--2.1]{hartman2002ode} applies and yields a unique $q\in C([0,T]\times[0,1];\mathbb R^2)$; inserting $q$ as a forcing into \eqref{eq:adjoint_p} and invoking the same result gives a unique $p\in C([0,T]\times[0,1];\mathbb R^2)$.
\end{proof}

The terminal condition \eqref{eq:adjoint_tc} reflects that the cost depends on $X(T,\cdot)$ only through $a(T,\cdot)$ via the regularized length. With $(p,q)$ at hand, we can now identify the Fr\'echet derivative of the reduced cost.

\begin{prop}[Reduced gradient via the adjoint system]\label{prop:adjoint}
Under the hypotheses of Theorem~\ref{prop:H3} and set $j_{\lambda,\varepsilon}(u):=\mathcal J_{\lambda,\varepsilon}(\psi)$. Then $j_{\lambda,\varepsilon}$ is Fr\'echet differentiable on $L^2(0,T;\mathbb R^N)$, and for every $\delta u\in L^2(0,T;\mathbb R^N)$ with $\eta:=\sum_{k=1}^N\delta u_k h_k$,
\begin{align}\label{eq:DJ}
Dj_{\lambda,\varepsilon}(u)[\delta u]
=\lambda\!\int_0^T\!\!(Mu(t))\cdot\delta u(t)\,dt
+\!\int_0^T\!\!\int_0^1\!\Bigl[p\cdot J\nabla_x\eta(t,X)+q\cdot J\nabla_x^2\eta(t,X)\,a\Bigr]ds\,dt,
\end{align}
or, componentwise, with $k=1,\dots,N$,
\begin{equation}\label{eq:grad_uk}
\frac{\partial j_{\lambda,\varepsilon}}{\partial u_k}(t)
=\lambda\,(Mu(t))_k+\int_0^1\!\Bigl[p(t,s)\cdot J\nabla h_k(X)+q(t,s)\cdot J\nabla_x^2 h_k(X)\,a(t,s)\Bigr]ds.
\end{equation}
\end{prop}

\begin{proof}
\textit{Step 1. Linearized state.}
Since each $h_k\in C^\infty(\overline\Omega)$, the right-hand side of \eqref{eq:state} is $C^\infty$ in $(X,a)$ and affine in $u$. Gr\"onwall's inequality applied to the difference of two state trajectories, with the $L^2$-integrable Lipschitz constants furnished by the pointwise bound \eqref{eq:pt_bound} and \eqref{eq:Ma}, yields
\[
\|(X_{u+\delta u}-X_u,\,a_{u+\delta u}-a_u)\|_{C([0,T]\times[0,1])}\le C\,\|\delta u\|_{L^2(0,T;\mathbb R^N)},
\]
and a second Gr\"onwall pass on the quadratic remainder gives an $O(\|\delta u\|_{L^2}^2)$ bound. Hence $u\mapsto(X,a)$ is Fr\'echet differentiable from $L^2(0,T;\mathbb R^N)$ into $C([0,T]\times[0,1];\mathbb R^2\times\mathbb R^2)$, and its derivative $(\delta X,\delta a):=(DX[\delta u],Da[\delta u])$ solves
\begin{subequations}
\begin{align}
\partial_t\delta X &= J\nabla_x^2\psi(t,X)\,\delta X+\sum_{k=1}^N\delta u_k(t)\,J\nabla h_k(X),\label{eq:linearized_X}\\
\partial_t\delta a &= J\nabla_x^2\psi(t,X)\,\delta a+J\nabla_x^3\psi(t,X)[\delta X,a]+\sum_{k=1}^N\delta u_k(t)\,J\nabla_x^2 h_k(X)\,a,\label{eq:linearized_a}\\
\delta X(0,s)&=0,\qquad \delta a(0,s)=0.\label{eq:linearized_ic}
\end{align}
\end{subequations}

\vskip 5pt
\noindent\textit{Step 2. Lagrangian identity.}
Pair \eqref{eq:linearized_X} with $p$ and \eqref{eq:linearized_a} with $q$ in $L^2((0,T)\times(0,1))$ and integrate by parts in $t$; the boundary conditions \eqref{eq:linearized_ic} and $p(T,\cdot)=0$ of \eqref{eq:adjoint_tc} give
\begin{equation}\label{eq:IBP}
\begin{aligned}
\int_0^T\!\!\int_0^1 p\cdot\partial_t\delta X\,ds\,dt &= -\int_0^T\!\!\int_0^1 \partial_t p\cdot\delta X\,ds\,dt,\\
\int_0^T\!\!\int_0^1 q\cdot\partial_t\delta a\,ds\,dt &= \int_0^1 q(T,s)\cdot\delta a(T,s)\,ds-\int_0^T\!\!\int_0^1 \partial_t q\cdot\delta a\,ds\,dt.
\end{aligned}
\end{equation}
Substituting \eqref{eq:linearized_X}--\eqref{eq:linearized_a} on the left and \eqref{eq:adjoint_p}--\eqref{eq:adjoint_q} on the right of \eqref{eq:IBP}, the transpose identities
\[
p\cdot J\nabla_x^2\psi\,\delta X=(J\nabla_x^2\psi)^\top p\cdot\delta X,
\qquad
q\cdot J\nabla_x^3\psi[\delta X,a]=\bigl(J\nabla_x^3\psi[\,\cdot\,,a]\bigr)^\top q\cdot\delta X,
\]
and the analogous identity for $q\cdot J\nabla_x^2\psi\,\delta a$, pair up the bulk contributions and cancel them. What remains is
\begin{align}\label{eq:Lagrangian_id}
\int_0^1 q(T,s)\cdot\delta a(T,s)\,ds
=\int_0^T\!\!\int_0^1\sum_{k=1}^N\delta u_k(t)\Bigl[p\cdot J\nabla h_k(X)+q\cdot J\nabla_x^2 h_k(X)\,a\Bigr]ds\,dt.
\end{align}

\vskip 5pt
\noindent\textit{Step 3. Differentiating the cost.}
By \eqref{eq:cost} and \eqref{eq:gram_identity}, $$\mathcal J_{\lambda,\varepsilon}(\psi)=\frac\lambda2\int_0^T u^\top Mu\,dt-\int_0^1\ell_\varepsilon(a(T,s))\,ds.$$ Differentiating in $\delta u$ and using $q(T,s)=-\nabla\ell_\varepsilon(a(T,s))$,
\[
Dj_{\lambda,\varepsilon}(u)[\delta u]
=\lambda\!\int_0^T\!\!(Mu)\cdot\delta u\,dt+\int_0^1 q(T,s)\cdot\delta a(T,s)\,ds,
\]
which combined with \eqref{eq:Lagrangian_id} yields \eqref{eq:DJ}. Taking $\delta u_\ell(t)=\delta_{k\ell}\varphi(t)$ for arbitrary $\varphi\in L^2(0,T)$ specializes \eqref{eq:DJ} to \eqref{eq:grad_uk}.
\end{proof}

Before turning to the numerical scheme, we record a consistency property linking the continuous gradient \eqref{eq:grad_uk} to its discrete counterpart introduced in Section~\ref{sec:numerics}.

\begin{remark}[Discrete gradient consistency]\label{rem:disc_cont_consistency}
The continuous penalty $$\frac\lambda 2\int_0^T u(t)^\top M u(t)\,dt,$$ equivalently $\frac\lambda 2\int_0^T\|\psi(t,\cdot)\|_{L^2(\Omega)}^2\,dt$ by \eqref{eq:gram_identity}, discretizes consistently to the trapezoidal sum $\frac{\Delta t}{2}\sum_n u_n^\top \widetilde M u_n$, with $\widetilde M$ a consistent discretization of $M$. In the implementation \eqref{eq:Jh_impl}--\eqref{eq:grad_impl} we work with a basis whose $L^2(\Omega)$ Gram matrix is diagonal (cf.\ Section~\ref{sec:ansatz}), so $\widetilde M=\mathrm{diag}(\gamma_k)$ with $\gamma_k=\|h_k\|_{L^2(\Omega)}^2$; in particular, the canonical case $\gamma_k\equiv 1$ corresponds to an $L^2(\Omega)$-orthonormal basis. Correspondingly, the $\lambda\bigl(Mu(t)\bigr)_k$ contribution in \eqref{eq:grad_uk} discretizes to $\Delta t\,\gamma_k u_{k,n}$, matching \eqref{eq:grad_impl}. Algebraic exactness of this discrete gradient is established in Remark~\ref{rem:discrete_exact}.
\end{remark}

For use in the numerical implementation of Section~\ref{sec:numerics}, we record that on $\mathcal U_{\mathrm{ad}}^{R_u}$ the velocity field and its gradient read
\[
\mathbf{v}_N(t,x)
=
\sum_{k=1}^N u_k(t)\,J\nabla h_k(x),
\qquad
\nabla_x\mathbf{v}_N(t,x)
=
\sum_{k=1}^N u_k(t)\,J\nabla_x^2 h_k(x).
\]
The coupled system \eqref{eq:state}--\eqref{eq:adjoint}--\eqref{eq:grad_uk} then constitutes the first-order necessary optimality conditions for the reduced optimal control problem.


\section{Numerical framework for the optimal control problem}\label{sec:numerics}

We now describe the full numerical pipeline for solving the optimal control problem \eqref{eq:ocp}: the control parameterization, the symplectic time integrator, the discrete adjoint and gradient computation, and the optimization algorithm.

\subsection{Finite-dimensional ansatz and control parameterization}\label{sec:ansatz}

Following the framework of Section~\ref{sec:ocp}, the stream function is expanded in a finite-dimensional basis $\{h_k\}_{k=1}^N\subset C^\infty(\overline\Omega)\cap H_0^1(\Omega)$:
\begin{equation*}
\psi_N(t,x)=\sum_{k=1}^N u_k(t)\,h_k(x),
\end{equation*}
where $u=(u_1,\dots,u_N)\in L^2(0,T;\mathbb R^N)$ are the control variables. The associated velocity fields $\mathbf v_k:=\nabla^\perp h_k=J\nabla h_k$ are divergence-free by construction. The specific choice of basis is problem-dependent; in Section~\ref{sec:optimal_flows}, we employ the cellular-flow basis, and in Section~\ref{sec:optimal_Doswell}, we use a Doswell-type multi-vortex velocity field.

\subsection{Symplectic ODE solver for the interface}\label{sec:symplectic}

The initial material interface $\Sigma(0)$ is discretized by $N_p$ marker points $X_0^1,\dots,X_0^{N_p}$, distributed uniformly with respect to arc length along $\Sigma(0)$.

Since the flow map \eqref{eq:hamiltonian_flow} is Hamiltonian, we employ a symplectic time integrator \cite{HLW06} to preserve the volume-preserving structure and prevent secular energy drift. We use the \emph{implicit midpoint rule}, the simplest symplectic Runge--Kutta method \cite{HLW06}: for $\dot y=f(t,y)$,
\begin{equation}\label{eq:midpoint_def}
y_{n+1}=y_n+\Delta t\,f\!\left(\frac{t_n+t_{n+1}}{2},\,\frac{y_n+y_{n+1}}{2}\right).
\end{equation}
This method is second-order accurate, A-stable, symmetric (self-adjoint), and symplectic. Because of its symmetry, transposing the linearized forward map yields the same midpoint scheme applied backward in time, so the discrete forward and adjoint solves are structurally consistent.

Applied to our setting, the marker update reads
\begin{equation}\label{eq:midpoint_marker}
X_{n+1}^j
=
X_n^j
+\Delta t\sum_{k=1}^N u_{k,n}\,
\mathbf v_k\!\left(\frac{X_n^j+X_{n+1}^j}{2}\right),
\end{equation}
where $u_{k,n}:=u_k(t_n)$ and $\mathbf v_k:=J\nabla h_k$. Throughout, the controls $u_k$ are taken to be piecewise constant on each interval $[t_n,t_{n+1})$, so that the velocity field is time-independent within each step and the implicit midpoint rule~\eqref{eq:midpoint_def} applies directly. The implicit equation is solved by five fixed-point iterations, initialized with an explicit Euler predictor.

The terminal polyline length is computed using the same $\varepsilon$-regularization as in the continuous cost \eqref{eq:cost}:
\begin{equation}\label{eq:polyline_length}
\mathcal L_{h,\varepsilon}(X_M)
:=
\sum_{j=1}^{N_p-1}\ell_\varepsilon(X_M^{j+1}-X_M^j),
\qquad
\ell_\varepsilon(z):=\sqrt{|z|^2+\varepsilon^2},
\end{equation}
where $M$ is the number of time steps and $X_M^j$ denotes the position of the $j$-th marker at the final time $t_M=T$. This quantity is the discrete geometric counterpart of the continuous length \eqref{eq:length}, and the smoothing parameter $\varepsilon>0$ ensures differentiability at zero chord length; the smoothing is a deliberate choice that aligns the cost with first-order conjugate-gradient methods at minimal practical cost (see Section~\ref{sec:numerics_exp}), with nonsmooth alternatives (ADMM \cite{Glowinski1992}, proximal-gradient methods on the unsmoothed length functional) left as future work. As $N_p\to\infty$ with the marker spacing $|X_M^{j+1}-X_M^j|=O(1/N_p)$, the polyline length $\mathcal L_{h,\varepsilon}(X_M)$ converges to the regularized continuous length $\int_0^1\ell_\varepsilon(\partial_sX(T,s))\,ds$ whenever the interface $s\mapsto X(T,s)$ is $W^{1,1}$-regular, and the gradient of $\mathcal L_{h,\varepsilon}$ with respect to marker positions approximates the shape derivative induced by the continuous tangent-based formulation. For interfaces in $W^{2,\infty}([0,1];\overline\Omega)$ with uniformly spaced markers $s_j=(j-1)/(N_p-1)$, a standard piecewise-linear quadrature estimate yields the rate
\begin{equation*}
\Bigl|\,\mathcal L_{h,\varepsilon}(X_M)-\int_0^1\ell_\varepsilon\!\bigl(\partial_sX(T,s)\bigr)\,ds\,\Bigr|
\le
\frac{C}{N_p}\,\|\partial_s^2X(T,\cdot)\|_{L^\infty([0,1])},
\end{equation*}
with $C$ depending only on $\varepsilon$ and $\|\partial_sX(T,\cdot)\|_{L^\infty}$, in line with Theorem~\ref{prop:H3}(ii) Step~3. The fully discrete adjoint constructed below is therefore consistent, in the marker-refinement limit, with the continuous $(p,q)$ sensitivity of Theorem~\ref{prop:H3}; in finite $N_p$ it is the algebraically exact adjoint of the implemented polyline objective.


\subsection{Discrete adjoint and gradient computation}\label{sec:discrete_adjoint}

We minimize the fully discrete objective
\begin{equation}\label{eq:Jh_impl}
J_h(u)
:=
-\mathcal L_{h,\varepsilon}(X_M)
+\frac{\Delta t}{2}\sum_{n=0}^{M-1}\sum_{k=1}^N
\gamma_k\,u_{k,n}^2,
\end{equation}
where $\gamma_k>0$ are the regularization weights and $u=\{u_{k,n}\}_{k=1,\dots,N}^{n=0,\dots,M-1}$ is the full set of discrete controls. This is the direct discretization of the continuous $L^2$ penalty $\frac\lambda2\int_0^T|u(t)|^2\,dt$ in \eqref{eq:cost} (with $\lambda$ absorbed into the per-component weights $\gamma_k$); the choice $\gamma_k\equiv\lambda$ recovers the canonical case, and nonuniform $\gamma_k$ allows per-component scaling without affecting the exactness of the discrete gradient derived below.

We emphasize that the adjoint construction below is the exact adjoint of the fully discrete marker/polyline objective \eqref{eq:Jh_impl}, obtained by transposing the linearized forward map of the implicit-midpoint marker dynamics. It is \emph{not} a direct discretization of the continuous $(X,a,p,q)$ adjoint system \eqref{eq:adjoint}: the continuous adjoint tracks separate adjoint variables $p$ and $q$ for position and tangent, whereas at the discrete level the tangent information is implicitly encoded in the finite differences $X_M^{j+1}-X_M^j$, so the discrete adjoint operates directly on marker positions. Both adjoints are adjoints of the same physical sensitivity at their respective levels (continuous and fully discrete), but the implemented algorithm is built from the discrete one and inherits its algebraic exactness. The terminal condition is the gradient of the regularized polyline length \eqref{eq:polyline_length}: Let $p_n^j\in\mathbb R^2$ denote the discrete adjoint variable at time step $n$ and marker $j$. Then
\[
p_M^j=-\partial_{X_M^j}\mathcal L_{h,\varepsilon}(X_M).
\]
Defining the regularized unit tangent vectors
\[
\tau_M^j:=\frac{X_M^{j+1}-X_M^j}{\sqrt{|X_M^{j+1}-X_M^j|^2+\varepsilon^2}},
\qquad j=1,\dots,N_p-1,
\]
the terminal adjoint values are
\begin{equation*}
p_M^1=\tau_M^1,
\qquad
p_M^{N_p}=-\tau_M^{N_p-1},
\qquad
p_M^j=\tau_M^j-\tau_M^{j-1}
\quad (2\le j\le N_p-1).
\end{equation*}

The backward adjoint step is also midpoint-based. Writing
\begin{equation}\label{eq:Gmid_impl}
G_n^j
:=
\sum_{k=1}^N u_{k,n}\,\nabla \mathbf v_k\!\left(\frac{X_n^j+X_{n+1}^j}{2}\right),
\end{equation}
the discrete adjoint update is
\begin{equation}\label{eq:adjoint_midpoint_impl}
p_n^j
=
p_{n+1}^j
+\Delta t\,(G_n^j)^\top\,\frac{p_n^j+p_{n+1}^j}{2}.
\end{equation}
This implicit equation is again solved by five fixed-point iterations, mirroring the forward solver.

The gradient of the discrete objective with respect to the controls is then assembled as
\begin{equation}\label{eq:grad_impl}
\frac{\partial J_h}{\partial u_{k,n}}
=
\gamma_k\,\Delta t\,u_{k,n}
+\Delta t\sum_{j=1}^{N_p}
\frac{p_n^j+p_{n+1}^j}{2}\cdot
\mathbf v_k\!\left(\frac{X_n^j+X_{n+1}^j}{2}\right),
\qquad k=1,\dots,N.
\end{equation}
\begin{remark}[Discrete gradient exactness]\label{rem:discrete_exact}
We verify algebraically that \eqref{eq:adjoint_midpoint_impl}--\eqref{eq:grad_impl} yield the exact gradient of $J_h$.
Consider the one-step forward map $F_n:X_n\mapsto X_{n+1}$ defined implicitly by~\eqref{eq:midpoint_marker}. When the fixed-point iteration converges, the implicit relation
\[
X_{n+1}=X_n+\Delta t\,f\!\Big(\frac{X_n+X_{n+1}}{2},u_n\Big)
\]
can be differentiated with respect to $X_n$, yielding $\Phi_n=I+\frac{\Delta t}{2}G_n(I+\Phi_n)$, hence
\[
\Big(I-\frac{\Delta t}{2}\,G_n\Big)\,\Phi_n
=I+\frac{\Delta t}{2}\,G_n,
\]
where $\Phi_n:=\partial X_{n+1}/\partial X_n$ is the one-step Jacobian and $G_n$ is the velocity-gradient matrix~\eqref{eq:Gmid_impl} evaluated at the midpoint. By the chain rule, the gradient of $J_h$ satisfies the backward recursion $p_n=\Phi_n^\top p_{n+1}$. Transposing the relation above gives
\[
\Phi_n^\top\Big(I-\frac{\Delta t}{2}\,G_n^\top\Big)
=I+\frac{\Delta t}{2}\,G_n^\top,
\]
so the backward recursion becomes
\[
\Big(I-\frac{\Delta t}{2}\,G_n^\top\Big)p_n
=\Big(I+\frac{\Delta t}{2}\,G_n^\top\Big)p_{n+1},
\]
which is precisely the midpoint adjoint step~\eqref{eq:adjoint_midpoint_impl}. Similarly, differentiating the forward map with respect to $u_{k,n}$ and contracting with $\frac{1}{2}(p_n+p_{n+1})$ gives the gradient formula~\eqref{eq:grad_impl}. This ``discretize-then-optimize'' consistency holds because the implicit midpoint rule is self-adjoint: transposing the linearized forward map yields the same midpoint scheme applied backward in time.
\end{remark}

In practice, five fixed-point iterations suffice for convergence to machine precision at the step sizes used in our experiments.

\subsection{Optimization algorithm}\label{sec:algorithm}

The control update is performed by a Polak--Ribi\`ere (PR) conjugate-gradient (CG) method with a \emph{projected} Armijo backtracking line search, combined with the closed-form projection onto the $L^2(0,T;\mathbb R^N)$-ball of radius $R_u$ that enforces the admissibility constraint $\|u\|_{L^2(0,T;\mathbb R^N)}\le R_u$ built into $\mathcal U_{\mathrm{ad}}^{R_u}$ in \eqref{eq:Uad}:
\begin{equation}\label{eq:projection}
\Pi_{B_{R_u}}(u):=\min\!\Bigl(1,\frac{R_u}{\|u\|_{L^2(0,T;\mathbb R^N)}}\Bigr)\,u.
\end{equation}
The Armijo decrease is checked at the projected trial point so that every accepted iterate is admissible by construction; the scheme reduces to the unconstrained PR-CG method whenever the ball is inactive along the line search. The complete procedure is stated in Algorithm~\ref{alg:optimize}.


\begin{algorithm}[t]
\caption{Adjoint-based optimization of interface length}\label{alg:optimize}
\begin{algorithmic}[1]
    \State \textbf{Input:} Initial interface $\gamma_0$, basis functions $\{h_k\}_{k=1}^N$, parameters $T,\lambda>0$, regularization weights $\{\gamma_k\}_{k=1}^N$, ball radius $R_u>0$, number of markers $N_p$, time steps $M$, Armijo constant $c_1>0$, stopping tolerance $\varepsilon_{\mathrm{tol}}>0$.
    \State \textbf{Output:} Optimized discrete controls $\{u_{k,n}^*\}$.
    \State \textbf{Initialize:} Set initial controls $\{u_{k,n}^{(0)}\}$, iteration counter $\ell:=0$.
    \State Compute $X_{\mathrm{traj}}^{(0)}$ by forward solve \eqref{eq:midpoint_marker} and $J_h^{(0)}:=J_h(u^{(0)})$.
    \State Compute discrete adjoint $\{p_n^{j,(0)}\}$ by backward solve \eqref{eq:adjoint_midpoint_impl}.
    \State Compute gradient $g^{(0)}$ via \eqref{eq:grad_impl}; set search direction $d^{(0)}:=-g^{(0)}$.
    \Repeat
        \State \textbf{(1) Projected line search.} Find step size $\eta>0$ satisfying the Armijo condition evaluated at the projected iterate:
        \[
            J_h\!\bigl(\Pi_{B_{R_u}}(u^{(\ell)}+\eta\,d^{(\ell)})\bigr)\le J_h(u^{(\ell)})+c_1\,\eta\,\langle g^{(\ell)},d^{(\ell)}\rangle,
        \]
        where $\Pi_{B_{R_u}}$ is given by \eqref{eq:projection}.
        \State \textbf{(2) Update controls.} Set $u^{(\ell+1)}:=\Pi_{B_{R_u}}(u^{(\ell)}+\eta\,d^{(\ell)})$.
        \State \textbf{(3) Forward solve.} Compute $X_{\mathrm{traj}}^{(\ell+1)}$ and $J_h^{(\ell+1)}$.
        \State \textbf{(4) Adjoint solve.} Compute $\{p_n^{j,(\ell+1)}\}$ by backward solve.
        \State \textbf{(5) Gradient and CG update.}
        \begin{enumerate}[(i)]
            \item Compute $g^{(\ell+1)}$ via \eqref{eq:grad_impl}.
            \item Compute Polak--Ribi\`ere parameter:
            $\beta^{(\ell)}:=\max\!\Big(0,\;\frac{\langle g^{(\ell+1)},g^{(\ell+1)}-g^{(\ell)}\rangle}{\|g^{(\ell)}\|^2}\Big)$.
            \item Update search direction: $d^{(\ell+1)}:=-g^{(\ell+1)}+\beta^{(\ell)}d^{(\ell)}$.
            \item If $\langle g^{(\ell+1)},d^{(\ell+1)}\rangle\ge 0$, reset $d^{(\ell+1)}:=-g^{(\ell+1)}$.
        \end{enumerate}
        \State Set $\ell\leftarrow\ell+1$.
    \Until{$\displaystyle\frac{|J_h^{(\ell)}-J_h^{(\ell-1)}|}{1+|J_h^{(\ell)}|}<\varepsilon_{\mathrm{tol}}$}
\end{algorithmic}
\end{algorithm}


\section{Numerical experiments}\label{sec:numerics_exp}

In this section, we demonstrate the effectiveness of the optimal control approach and validate the proposed framework. We first analyze the mixing properties of two stationary velocity fields (Section~\ref{sec:stationary_flows}), then apply the optimization algorithm to design optimal time-dependent flows for the cellular-flow (Section~\ref{sec:optimal_flows}) and Doswell frontogenesis (Section~\ref{sec:optimal_Doswell}) settings. A systematic comparison with the Eulerian Sobolev-norm optimizer of \cite{hu2026structure} is presented in Section~\ref{sec:comparison}. We then investigate the influence of the number of basis functions by extending to $N=4$ modes (Section~\ref{sec:N4}). Throughout this section the smoothing parameter $\varepsilon$ in \eqref{eq:ell} is fixed at $10^{-8}$, several orders of magnitude below the typical chord length, so that the regularization is numerically inactive in all reported experiments. The code is implemented in Python using PyTorch and publicly available at \url{https://github.com/DCN-FAU-AvH/Mixing-Hamiltolnian}.

Throughout this section, the initial scalar distribution for the Eulerian transport solve is
\begin{equation}\label{eq:theta_init_1}
\theta^0(x_1,x_2) := \tanh\!\left(\frac{x_2-0.5}{0.01}\right),
\qquad (x_1,x_2)\in\Omega,
\end{equation}
and the initial material interface is the zero contour $\Sigma(0)=\{(x_1,x_2)\in\Omega:x_2=0.5\}$. In all optimization experiments, the regularization weights are set to $\gamma_1=\gamma_2=10^{-5}$ unless stated otherwise. All runs use the projected-gradient Algorithm~\ref{alg:optimize} with ball radius $R_u=10$; the terminal ratio $\|u^*\|_{L^2(0,T;\mathbb R^N)}/R_u$ is reported for each experiment to indicate whether the projection step is active (saturates the ball) or inactive.

\subsection{Stationary flows}\label{sec:stationary_flows}

\subsubsection{Cellular flow}
Cellular flows~\cite{Fannjiang2006Quenching} are standard benchmark velocity fields in the study of fluid mixing. We consider the steady divergence-free velocity field on $\Omega=(0,1)^2$:
\begin{equation}\label{eq:cellular}
\mathbf{v}(x_1,x_2)=
\begin{bmatrix}
-\pi\sin(\pi x_1)\cos(\pi x_2)\\[0.3em]
\ \ \pi\cos(\pi x_1)\sin(\pi x_2)
\end{bmatrix},
\end{equation}
corresponding to the stream function $h(x_1,x_2)=\sin(\pi x_1)\sin(\pi x_2)$.

We simulate over $t\in[0,5]$ with Eulerian mesh $\Delta x=\Delta t=1/1000$ and $N_p=100{,}000$ Lagrangian markers. Figures~\ref{fig:interface_cellular}--\ref{fig:theta_cellular} show the interface and scalar field evolution. As quantified in Figure~\ref{fig:cellular_mixnorm-length}, the interface length grows approximately linearly, consistent with polynomial decay of $\|\theta_h\|_{\dot{H}^{-1}(\Omega)}$. This algebraic behavior is consistent with the theoretical mixing literature for smooth autonomous two-dimensional Hamiltonian flows, in which the action--angle structure precludes genuinely exponential mix-norm decay even in the presence of local hyperbolic stretching; exponential mixing in this setting is associated instead with genuinely time-dependent mechanisms, such as alternating shears, random protocols, or uniformly hyperbolic maps \cite{zelati2024mixing}. In Sections~\ref{sec:optimal_flows}--\ref{sec:optimal_Doswell} we therefore expect the optimized \emph{time-dependent} controls, rather than any stationary refinement of the present field, to be the relevant route to faster decay; the rates reported there should accordingly be read as fitted exponential rates over the simulated horizon, not as asymptotic statements for arbitrary incompressible flows.

\begin{remark}[Local versus global sensitivity]\label{rem:local_global}
A consistent observation across all experiments is that the interface length begins to grow immediately, whereas the $\dot{H}^{-1}$ mix-norm remains nearly constant during an initial transient before decaying. This reflects the different nature of the two quantities: the interface length is a \emph{local} geometric measure that responds as soon as stretching occurs anywhere along the curve, while the $\dot{H}^{-1}$ norm is a \emph{global} measure that requires fine-scale filaments to affect the large-scale scalar distribution before registering decay \cite{mathew2005multiscale, thiffeault2012using}.
\end{remark}

\begin{figure}[ht]
  \centering
  \begin{subfigure}[b]{0.16\textwidth}
    \includegraphics[width=\textwidth]{./images/cellular/interface_t0.jpg}
    \caption{$t=0$}
  \end{subfigure}
  \begin{subfigure}[b]{0.16\textwidth}
    \includegraphics[width=\textwidth]{./images/cellular/interface_t1.jpg}
    \caption{$t=1$}
  \end{subfigure}
  \begin{subfigure}[b]{0.16\textwidth}
    \includegraphics[width=\textwidth]{./images/cellular/interface_t2.jpg}
    \caption{$t=2$}
  \end{subfigure}
  \begin{subfigure}[b]{0.16\textwidth}
    \includegraphics[width=\textwidth]{./images/cellular/interface_t3.jpg}
    \caption{$t=3$}
  \end{subfigure}
  \begin{subfigure}[b]{0.16\textwidth}
    \includegraphics[width=\textwidth]{./images/cellular/interface_t4.jpg}
    \caption{$t=4$}
  \end{subfigure}
  \begin{subfigure}[b]{0.16\textwidth}
    \includegraphics[width=\textwidth]{./images/cellular/interface_t5.jpg}
    \caption{$t=5$}
  \end{subfigure}
  \caption{Evolution of the interface under the steady cellular flow \eqref{eq:cellular} for $t\in[0,5]$.}
  \label{fig:interface_cellular}
\end{figure}

\begin{figure}[ht]
  \centering
  \begin{subfigure}[b]{0.16\textwidth}
    \includegraphics[width=\textwidth]{./images/cellular/theta_t0.jpg}
    \caption{$t=0$}
  \end{subfigure}
  \begin{subfigure}[b]{0.16\textwidth}
    \includegraphics[width=\textwidth]{./images/cellular/theta_t1.jpg}
    \caption{$t=1$}
  \end{subfigure}
  \begin{subfigure}[b]{0.16\textwidth}
    \includegraphics[width=\textwidth]{./images/cellular/theta_t2.jpg}
    \caption{$t=2$}
  \end{subfigure}
  \begin{subfigure}[b]{0.16\textwidth}
    \includegraphics[width=\textwidth]{./images/cellular/theta_t3.jpg}
    \caption{$t=3$}
  \end{subfigure}
  \begin{subfigure}[b]{0.16\textwidth}
    \includegraphics[width=\textwidth]{./images/cellular/theta_t4.jpg}
    \caption{$t=4$}
  \end{subfigure}
  \begin{subfigure}[b]{0.16\textwidth}
    \includegraphics[width=\textwidth]{./images/cellular/theta_t5.jpg}
    \caption{$t=5$}
  \end{subfigure}
  \caption{Evolution of $\theta_h$ under the steady cellular flow for $t\in[0,5]$.}
  \label{fig:theta_cellular}
\end{figure}

\begin{figure}[ht]
  \centering
  \begin{subfigure}[b]{0.475\textwidth}
    \centering
    \includegraphics[height=3cm]{./images/cellular/interface_length.png}
    \caption{Interface length $\mathcal{L}(t)$ (linear growth).}
    \label{fig:cellular_length}
  \end{subfigure}
   \begin{subfigure}[b]{0.49\textwidth}
    \centering
    \includegraphics[height=3cm]{./images/cellular/mixnorm.png}
    \caption{$\dot{H}^{-1}$ mix-norm (polynomial decay).}
    \label{fig:cellular_mixnorm}
  \end{subfigure}
    \caption{Interface growth and mix-norm decay under the steady cellular flow.}
  \label{fig:cellular_mixnorm-length}
\end{figure}



\subsubsection{Doswell frontogenesis} \label{sec:Doswell-stable}
To verify our observations in a setting with circular geometry, we consider the Doswell frontogenesis model \cite{doswell1984kinematic}. The domain is the disc $\Omega=\{\mathbf{x}\in\mathbb{R}^2:|\mathbf{x}-\mathbf{x}_c|<R\}$ with center $\mathbf{x}_c=(0.5,0.5)$ and radius $R=0.5$, so that the initial interface $\Sigma(0)=\{x_2=0.5\}$ passes through the disc center. The velocity field is
\begin{equation}
\label{eq:v_Doswell}
\mathbf{v}(x_1,x_2)=
 \begin{bmatrix}
      -(x_2-0.5)\,g(r)          \\[0.3em]
        \phantom{-}(x_1-0.5)\,g(r)
     \end{bmatrix},
\end{equation}
where $r:=\sqrt{(x_1-0.5)^2+(x_2-0.5)^2}$, $\overline{v}=2.59807$, and the angular velocity profile is
\[
    g(r):=\frac{\overline{v}}{r}\,\operatorname{sech}^2(r)\,\tanh(r).
\]
The velocity field \eqref{eq:v_Doswell} is divergence-free and tangent to $\partial\Omega$. The transport equation is discretized on a polar grid ($N_r=N_\varphi=1000$) over $t\in[0,10]$ with $\Delta t=1/1000$ and $N_p=100{,}000$ markers. As shown in Figures~\ref{fig:interface_Doswell}--\ref{fig:Doswell_mixnorm-length}, the same qualitative behavior is observed: linear interface growth and polynomial mix-norm decay.

\begin{figure}[ht]
  \centering
  \begin{subfigure}[b]{0.16\textwidth}
    \includegraphics[width=\textwidth]{./images/Doswell/interface_t0.jpg}
    \caption{$t=0$}
  \end{subfigure}
  \begin{subfigure}[b]{0.16\textwidth}
    \includegraphics[width=\textwidth]{./images/Doswell/interface_t1.jpg}
    \caption{$t=2$}
  \end{subfigure}
  \begin{subfigure}[b]{0.16\textwidth}
    \includegraphics[width=\textwidth]{./images/Doswell/interface_t2.jpg}
    \caption{$t=4$}
  \end{subfigure}
  \begin{subfigure}[b]{0.16\textwidth}
    \includegraphics[width=\textwidth]{./images/Doswell/interface_t3.jpg}
    \caption{$t=6$}
  \end{subfigure}
  \begin{subfigure}[b]{0.16\textwidth}
    \includegraphics[width=\textwidth]{./images/Doswell/interface_t4.jpg}
    \caption{$t=8$}
  \end{subfigure}
  \begin{subfigure}[b]{0.16\textwidth}
    \includegraphics[width=\textwidth]{./images/Doswell/interface_t5.jpg}
    \caption{$t=10$}
  \end{subfigure}
  \caption{Evolution of the interface under Doswell frontogenesis for $t\in[0,10]$.}
  \label{fig:interface_Doswell}
\end{figure}

\begin{figure}[ht]
  \centering
  \begin{subfigure}[b]{0.16\textwidth}
    \includegraphics[width=\textwidth]{./images/Doswell/theta_t0.jpg}
    \caption{$t=0$}
  \end{subfigure}
  \begin{subfigure}[b]{0.16\textwidth}
    \includegraphics[width=\textwidth]{./images/Doswell/theta_t1.jpg}
    \caption{$t=2$}
  \end{subfigure}
  \begin{subfigure}[b]{0.16\textwidth}
    \includegraphics[width=\textwidth]{./images/Doswell/theta_t2.jpg}
    \caption{$t=4$}
  \end{subfigure}
  \begin{subfigure}[b]{0.16\textwidth}
    \includegraphics[width=\textwidth]{./images/Doswell/theta_t3.jpg}
    \caption{$t=6$}
  \end{subfigure}
  \begin{subfigure}[b]{0.16\textwidth}
    \includegraphics[width=\textwidth]{./images/Doswell/theta_t4.jpg}
    \caption{$t=8$}
  \end{subfigure}
  \begin{subfigure}[b]{0.16\textwidth}
    \includegraphics[width=\textwidth]{./images/Doswell/theta_t5.jpg}
    \caption{$t=10$}
  \end{subfigure}
  \caption{Evolution of $\theta_h$ under Doswell frontogenesis for $t\in[0,10]$.}
  \label{fig:theta_Doswell}
\end{figure}

\begin{figure}[ht]
  \centering
  \begin{subfigure}[b]{0.475\textwidth}
    \centering
    \includegraphics[height=3cm]{./images/Doswell/interface_length.png}
    \caption{Interface length $\mathcal{L}(t)$ (linear growth).}
    \label{fig:Doswell_length}
  \end{subfigure}
  \begin{subfigure}[b]{0.49\textwidth}
    \centering
    \includegraphics[height=3cm]{./images/Doswell/mixnorm.png}
    \caption{$\dot{H}^{-1}$ mix-norm (polynomial decay).}
    \label{fig:Doswell_mixnorm}
  \end{subfigure}
    \caption{Interface growth and mix-norm decay under the Doswell velocity field.}
  \label{fig:Doswell_mixnorm-length}
\end{figure}




\subsection{Optimized cellular-flow mixers}\label{sec:optimal_flows}
Having established the limitations of stationary mixing, we now apply the adjoint-based optimization of Algorithm~\ref{alg:optimize} to design time-dependent stirring strategies. We restrict the stream function to $N=2$ modes with basis functions $h_k(x_1,x_2)=\sin(k\pi x_1)\sin(k\pi x_2)$, $k\in\{1,2\}$. The optimization is performed on $\Omega=(0,1)^2$ over the time horizon $t\in[0,1]$ with $M=1000$ time steps and $N_p=100{,}000$ markers. To evaluate the $\dot{H}^{-1}$ mix-norm, we use a finite-volume Eulerian transport solve with $\Delta x=1/1000$; the same resolution is used in all subsequent experiments.

\subsubsection{Constant initial guess}
We first initialize the controls as
\begin{equation}
\label{eq:control_cellular_1}
    u_1(t)=1,\quad u_2(t)=1, \quad t\in[0,1].
\end{equation}

Figure~\ref{fig:interface_Cellular-1} shows the geometric stretching of the interface under the optimized flow. Figure~\ref{fig:Cellular_1_mixnorm-length} shows that the optimized velocity field produces exponential growth of the interface length at rate $5.74$, accompanied by an exponential decay of the $\dot{H}^{-1}$ mix-norm at rate $2.59$, in contrast to the polynomial decay observed in the stationary benchmarks. Figure~\ref{fig:theta_Cellular-1} shows the corresponding Eulerian scalar field, and the optimized controls $u_1(t)$ and $u_2(t)$ are displayed in Figure~\ref{fig:Cellular_1_u1u2}. At convergence the control attains $\|u^*\|_{L^2}/R_u = 0.16$, so the projection ball is inactive.

\begin{figure}[H]
  \centering
  \begin{subfigure}[b]{0.16\textwidth}
    \includegraphics[width=\textwidth]{./images/Cellular-1/interface_t0.jpg}
    \caption{$t=0$}
  \end{subfigure}
  \begin{subfigure}[b]{0.16\textwidth}
    \includegraphics[width=\textwidth]{./images/Cellular-1/interface_t1.jpg}
    \caption{$t=0.2$}
  \end{subfigure}
  \begin{subfigure}[b]{0.16\textwidth}
    \includegraphics[width=\textwidth]{./images/Cellular-1/interface_t2.jpg}
    \caption{$t=0.4$}
  \end{subfigure}
  \begin{subfigure}[b]{0.16\textwidth}
    \includegraphics[width=\textwidth]{./images/Cellular-1/interface_t3.jpg}
    \caption{$t=0.6$}
  \end{subfigure}
  \begin{subfigure}[b]{0.16\textwidth}
    \includegraphics[width=\textwidth]{./images/Cellular-1/interface_t4.jpg}
    \caption{$t=0.8$}
  \end{subfigure}
  \begin{subfigure}[b]{0.16\textwidth}
    \includegraphics[width=\textwidth]{./images/Cellular-1/interface_t5.jpg}
    \caption{$t=1.0$}
  \end{subfigure}
  \caption{Interface evolution under the optimized cellular flow (initial guess \eqref{eq:control_cellular_1}).}
  \label{fig:interface_Cellular-1}
\end{figure}

\begin{figure}[H]
  \centering
  \begin{subfigure}[b]{0.16\textwidth}
    \includegraphics[width=\textwidth]{./images/Cellular-1/theta_t0.jpg}
    \caption{$t=0$}
  \end{subfigure}
  \begin{subfigure}[b]{0.16\textwidth}
    \includegraphics[width=\textwidth]{./images/Cellular-1/theta_t1.jpg}
    \caption{$t=0.2$}
  \end{subfigure}
  \begin{subfigure}[b]{0.16\textwidth}
    \includegraphics[width=\textwidth]{./images/Cellular-1/theta_t2.jpg}
    \caption{$t=0.4$}
  \end{subfigure}
  \begin{subfigure}[b]{0.16\textwidth}
    \includegraphics[width=\textwidth]{./images/Cellular-1/theta_t3.jpg}
    \caption{$t=0.6$}
  \end{subfigure}
  \begin{subfigure}[b]{0.16\textwidth}
    \includegraphics[width=\textwidth]{./images/Cellular-1/theta_t4.jpg}
    \caption{$t=0.8$}
  \end{subfigure}
  \begin{subfigure}[b]{0.16\textwidth}
    \includegraphics[width=\textwidth]{./images/Cellular-1/theta_t5.jpg}
    \caption{$t=1.0$}
  \end{subfigure}
  \caption{Scalar field $\theta_h$ under the optimized cellular flow (initial guess \eqref{eq:control_cellular_1}).}
  \label{fig:theta_Cellular-1}
\end{figure}

\begin{figure}[H]
  \centering
  \begin{subfigure}[b]{0.32\textwidth}
    \includegraphics[width=\textwidth]{./images/Cellular-1/interface_length.png}
    \caption{Interface length $\mathcal{L}(t)$.}
    \label{fig:Cellular_1_length}
  \end{subfigure}
  \begin{subfigure}[b]{0.32\textwidth}
    \includegraphics[width=\textwidth]{./images/Cellular-1/mixnorm.png}
    \caption{$\dot{H}^{-1}$ mix-norm.}
    \label{fig:Cellular_1_mixnorm}
  \end{subfigure}
  \begin{subfigure}[b]{0.32\textwidth}
    \includegraphics[width=\textwidth]{./images/Cellular-1/u1u2.png}
    \caption{Optimal $u_1(t),u_2(t)$.}
    \label{fig:Cellular_1_u1u2}
  \end{subfigure}
    \caption{Quantitative results under the optimized cellular flow (initial guess \eqref{eq:control_cellular_1}).}
  \label{fig:Cellular_1_mixnorm-length}
\end{figure}



\subsubsection{Oscillatory initial guess}
To test robustness with respect to the initial guess, we perform a second optimization with the oscillatory initialization
\begin{equation}
\label{eq:control_cellular_2}
    u_1(t)=\cos(\pi t/2),\quad u_2(t)=\sin(\pi t/2), \quad t\in[0,1].
\end{equation}

Despite the different starting point, the optimized interface and scalar field evolution are qualitatively similar to the constant-initial-guess case (Figures~\ref{fig:interface_Cellular-1}--\ref{fig:theta_Cellular-1}). The interface grows exponentially at rate $5.12$ and the mix-norm decays exponentially at rate $2.28$ (Figure~\ref{fig:Cellular_2_mixnorm-length}), confirming that the algorithm is robust with respect to initialization. The terminal ratio is $\|u^*\|_{L^2}/R_u = 0.15$ (projection inactive).

\begin{figure}[ht]
  \centering
  \begin{subfigure}[b]{0.32\textwidth}
    \includegraphics[width=\textwidth]{./images/Cellular-2/interface_length.png}
    \caption{Interface length $\mathcal{L}(t)$.}
    \label{fig:Cellular_2_length}
  \end{subfigure}
  \begin{subfigure}[b]{0.34\textwidth}
    \includegraphics[width=\textwidth]{./images/Cellular-2/mixnorm.png}
    \caption{$\dot{H}^{-1}$ mix-norm.}
    \label{fig:Cellular_2_mixnorm}
  \end{subfigure}
  \begin{subfigure}[b]{0.32\textwidth}
    \includegraphics[width=\textwidth]{./images/Cellular-2/u1u2.png}
    \caption{Optimal $u_1(t),u_2(t)$.}
    \label{fig:Cellular_2_u1u2}
  \end{subfigure}
    \caption{Quantitative results under the optimized cellular flow (initial guess \eqref{eq:control_cellular_2}).}
  \label{fig:Cellular_2_mixnorm-length}
\end{figure}




\subsection{Optimized Doswell-type mixers}\label{sec:optimal_Doswell}
We next consider a more complex flow configuration on a circular domain to demonstrate the flexibility of our framework beyond the Cartesian setting. The domain is the same disc $\Omega = \{ \mathbf{x} \in \mathbb{R}^2 : |\mathbf{x} - \mathbf{x}_c| < R \}$ with center $\mathbf{x}_c = (0.5, 0.5)$ and radius $R=0.5$ as in the stationary benchmark of Section~\ref{sec:Doswell-stable}. The optimization is performed over $t\in[0,5]$ with $M=5000$ time steps and $N_p=100{,}000$ markers. For the Eulerian transport solve used in mix-norm evaluation, we use a polar grid with $N_r=1000$ radial cells and $N_\varphi=1000$ angular cells.

The velocity field is parameterized by $N=2$ basis flows. The first, $\mathbf{b}_1$, is the standard Doswell frontogenesis field \eqref{eq:v_Doswell} centered at $\mathbf{x}_c$. The second, $\mathbf{b}_2$, is a five-cell (multi-vortex) Doswell-type field formed by superposing five Doswell vortices supported on five embedded discs within $\Omega$: four congruent discs arranged around the periphery and one centered at $\mathbf{x}_c$. To enforce no-penetration at each disc boundary and to obtain $C^1$ spatial regularity when the field is extended by zero outside the disc, each vortex is multiplied by the radial cutoff
\[
C(r)=\bigl(1-(r/R_c)^2\bigr)^3 \quad (0\le r<R_c),
\qquad C(r)=0 \quad (r\ge R_c),
\]
where $R_c$ denotes the disc radius and $r$ is the distance to the corresponding vortex center. This choice satisfies $C(R_c)=0$ and $C'(R_c)=0$. The resulting field $\mathbf{b}_2$ is incompressible in $\Omega$ and satisfies $\mathbf{b}_2\cdot\mathbf{n}=0$ on $\partial\Omega$. The cubic cutoff only guarantees $C^1$ regularity across $\{r=R_c\}$, so the associated stream function is not in $C^\infty(\overline\Omega)$ and the basis $\{\mathbf{b}_1,\mathbf{b}_2\}$ falls outside the smoothness hypothesis of Theorem~\ref{prop:H3}. The experiments below should therefore be understood as a numerical demonstration on a physically motivated basis, rather than as an illustration of the existence theorem.

\subsubsection{Constant initial guess}
We initialize the controls as in \eqref{eq:control_cellular_1}. Figures~\ref{fig:interface_Doswell-1} and \ref{fig:theta_Doswell-1} display the optimized interface and scalar field evolution, respectively. Despite the more complex flow geometry, the optimizer again produces chaotic advection: the interface length grows exponentially at rate $0.92$ (Figure~\ref{fig:Doswell_1_length}), driving an exponential decay of the mix-norm at rate $0.33$ (Figure~\ref{fig:Doswell_1_mixnorm}). The optimized controls are shown in Figure~\ref{fig:Doswell_1_u1u2}. At convergence the control attains $\|u^*\|_{L^2}/R_u = 1.00$, i.e., the iterate sits on the boundary of the projection ball; this is the single experiment in which the projection step of Algorithm~\ref{alg:optimize} is active.

\begin{figure}[ht]
  \centering
  \begin{subfigure}[b]{0.16\textwidth}
    \includegraphics[width=\textwidth]{./images/Doswell-1/interface_t0.jpg}
    \caption{$t=0$}
  \end{subfigure}
  \begin{subfigure}[b]{0.16\textwidth}
    \includegraphics[width=\textwidth]{./images/Doswell-1/interface_t1.jpg}
    \caption{$t=1$}
  \end{subfigure}
  \begin{subfigure}[b]{0.16\textwidth}
    \includegraphics[width=\textwidth]{./images/Doswell-1/interface_t2.jpg}
    \caption{$t=2$}
  \end{subfigure}
  \begin{subfigure}[b]{0.16\textwidth}
    \includegraphics[width=\textwidth]{./images/Doswell-1/interface_t3.jpg}
    \caption{$t=3$}
  \end{subfigure}
  \begin{subfigure}[b]{0.16\textwidth}
    \includegraphics[width=\textwidth]{./images/Doswell-1/interface_t4.jpg}
    \caption{$t=4$}
  \end{subfigure}
  \begin{subfigure}[b]{0.16\textwidth}
    \includegraphics[width=\textwidth]{./images/Doswell-1/interface_t5.jpg}
    \caption{$t=5$}
  \end{subfigure}
  \caption{Interface evolution under the optimized Doswell flow (initial guess \eqref{eq:control_cellular_1}).}
  \label{fig:interface_Doswell-1}
\end{figure}

\begin{figure}[ht]
  \centering
  \begin{subfigure}[b]{0.16\textwidth}
    \includegraphics[width=\textwidth]{./images/Doswell-1/theta_t0.jpg}
    \caption{$t=0$}
  \end{subfigure}
  \begin{subfigure}[b]{0.16\textwidth}
    \includegraphics[width=\textwidth]{./images/Doswell-1/theta_t1.jpg}
    \caption{$t=1$}
  \end{subfigure}
  \begin{subfigure}[b]{0.16\textwidth}
    \includegraphics[width=\textwidth]{./images/Doswell-1/theta_t2.jpg}
    \caption{$t=2$}
  \end{subfigure}
  \begin{subfigure}[b]{0.16\textwidth}
    \includegraphics[width=\textwidth]{./images/Doswell-1/theta_t3.jpg}
    \caption{$t=3$}
  \end{subfigure}
  \begin{subfigure}[b]{0.16\textwidth}
    \includegraphics[width=\textwidth]{./images/Doswell-1/theta_t4.jpg}
    \caption{$t=4$}
  \end{subfigure}
  \begin{subfigure}[b]{0.16\textwidth}
    \includegraphics[width=\textwidth]{./images/Doswell-1/theta_t5.jpg}
    \caption{$t=5$}
  \end{subfigure}
  \caption{Scalar field $\theta_h$ under the optimized Doswell flow (initial guess \eqref{eq:control_cellular_1}).}
  \label{fig:theta_Doswell-1}
\end{figure}

\begin{figure}[ht]
  \centering
  \begin{subfigure}[b]{0.32\textwidth}
    \includegraphics[width=\textwidth]{./images/Doswell-1/interface_length.png}
    \caption{Interface length $\mathcal{L}(t)$.}
    \label{fig:Doswell_1_length}
  \end{subfigure}
  \begin{subfigure}[b]{0.34\textwidth}
    \includegraphics[width=\textwidth]{./images/Doswell-1/mixnorm.png}
    \caption{$\dot{H}^{-1}$ mix-norm.}
    \label{fig:Doswell_1_mixnorm}
  \end{subfigure}
  \begin{subfigure}[b]{0.32\textwidth}
    \includegraphics[width=\textwidth]{./images/Doswell-1/u1u2.png}
    \caption{Optimal $u_1(t),u_2(t)$.}
    \label{fig:Doswell_1_u1u2}
  \end{subfigure}
    \caption{Quantitative results under optimized Doswell flow (initial guess \eqref{eq:control_cellular_1}).}
  \label{fig:Doswell_1_mixnorm-length}
\end{figure}






\subsubsection{Oscillatory initial guess}
Starting from the oscillatory initialization \eqref{eq:control_cellular_2}, the optimizer converges to a qualitatively similar solution, with interface and scalar field evolution close to the constant-initial-guess case (Figures~\ref{fig:interface_Doswell-1}--\ref{fig:theta_Doswell-1}). The interface length grows exponentially at rate $0.70$ and the mix-norm decays at rate $0.41$ (Figure~\ref{fig:Doswell_2_mixnorm-length}), again showing robustness with respect to initialization. The terminal ratio is $\|u^*\|_{L^2}/R_u = 0.70$ (projection inactive, but closer to saturation than the cellular cases).

\begin{figure}[ht]
  \centering
  \begin{subfigure}[b]{0.32\textwidth}
    \includegraphics[width=\textwidth]{./images/Doswell-2/interface_length.png}
    \caption{Interface length $\mathcal{L}(t)$.}
    \label{fig:Doswell_2_length}
  \end{subfigure}
  \begin{subfigure}[b]{0.34\textwidth}
    \includegraphics[width=\textwidth]{./images/Doswell-2/mixnorm.png}
    \caption{$\dot{H}^{-1}$ mix-norm.}
    \label{fig:Doswell_2_mixnorm}
  \end{subfigure}
  \begin{subfigure}[b]{0.32\textwidth}
    \includegraphics[width=\textwidth]{./images/Doswell-2/u1u2.png}
    \caption{Optimal $u_1(t),u_2(t)$.}
    \label{fig:Doswell_2_u1u2}
  \end{subfigure}
    \caption{Quantitative results under optimized Doswell flow (initial guess \eqref{eq:control_cellular_2}).}
  \label{fig:Doswell_2_mixnorm-length}
\end{figure}



Across the four optimization experiments, the algorithm transforms the polynomial mixing of the stationary benchmarks into rates well fitted by exponentials over the simulated horizon~\cite{zheng2023numerical, hu2026structure}. The Doswell-optimized cases yield lower $\dot{H}^{-1}$ rates ($0.33$, $0.41$) than the cellular ones ($2.59$, $2.28$) despite producing long, geometrically complex spirals: the Doswell mechanism coils the interface within the disc's coherent circular geometry, so length grows quickly while the global $\dot{H}^{-1}$ cancellations improve only moderately. This decoupling, consistent with Remark~\ref{rem:critical_mixnorm}, reinforces that all optimized controls must ultimately be assessed on the full transport equation.

\subsection{Comparison with the Sobolev-norm optimizer}\label{sec:comparison}

As discussions in Section~\ref{sec:problem}, the Hamiltonian interface model serves as a reduced-order surrogate for the original transport equation-constrained mixing problem. The natural validation of any such reduced model is to compare the performance of the control it produces against that of the control obtained by direct optimization on the full system, both evaluated on the same reference dynamics. We now carry out this comparison, using the Eulerian Sobolev-norm optimizer of \cite{hu2026structure} as the reference. That method minimizes $\frac{1}{2}\|\theta(T)\|_{\dot{H}^{-1}(\Omega)}^2+\frac{\gamma}{2}\int_0^T\|\mathbf{v}(t)\|^2\,dt$ via a structure-preserving finite-volume and Crank--Nicolson discretization.

The Sobolev-norm optimizer requires solving the full two-dimensional transport equation on an $N_{\mathrm{cell}}\times N_{\mathrm{cell}}$ grid (with $N_{\mathrm{cell}}=500$ as in \cite{hu2026structure}) at each forward step, together with a corresponding adjoint PDE solve of the same dimension, and additionally evaluating the $\dot{H}^{-1}$ norm via an elliptic solve. By contrast, our interface-length optimizer tracks only $N_p=100{,}000$ marker points along the one-dimensional material curve and computes the adjoint along the same trajectories. The per-iteration cost therefore scales as $O(N_p M)$ rather than $O(N_{\mathrm{cell}}^2 M)$, where $M$ is the number of time steps.

To ensure a fair comparison, the two optimizers run on the same finite-dimensional basis $(\mathbf b_1,\mathbf b_2)$, time horizon $T$, time step $\Delta t$, initial scalar~\eqref{eq:theta_init_1}, initial interface $\Sigma(0)$, constant initial control, $L^2$ penalty weights, projection onto $\|u\|_{L^2}\le R_u$, optimizer (Polak--Ribi\`ere CG with Armijo line search) and tolerance $\varepsilon_{\mathrm{tol}}=10^{-6}$; they differ only in the objective (interface length vs.\ $\dot{H}^{-1}$) and in the corresponding forward/adjoint solver. The optimized controls are evaluated on a common reference Eulerian solver ($N_{\mathrm{cell}}=1000$; $M_t=1000$ for cellular and $M_t=5000$ for Doswell), and Table~\ref{tab:optimizer_comparison} reports the resulting $\dot{H}^{-1}$ decay rate (least-squares regression of $\log\|\theta\|_{\dot{H}^{-1}}$ vs.\ $t$ over $[0,T]$, no burn-in) together with wall-clock time on the hardware described above.

\begin{table}[ht]
\centering
\begin{tabular}{|c|c|c|c|c|}
\hline
\multirow{2}{*}{Flow setting}
& \multicolumn{2}{c|}{Sobolev-norm optimizer}
& \multicolumn{2}{c|}{Interface-length optimizer} \\
\cline{2-5}
& Time (s) & Decay rate & Time (s) & Decay rate \\ \hline
Cellular with \eqref{eq:control_cellular_1} & 10522.89 & 2.52 & 1363.90 & 2.59 \\ \hline
Cellular with \eqref{eq:control_cellular_2} & 19423.52 & 1.81 & 894.07 & 2.28 \\ \hline
Doswell with \eqref{eq:control_cellular_1} & 18324.59 & 0.30 & 6714.30 & 0.33 \\ \hline
Doswell with \eqref{eq:control_cellular_2} & 23741.28 & 0.25 & 7634.20 & 0.41 \\ \hline
\end{tabular}
\caption{Comparison of the Sobolev-norm optimizer \cite{hu2026structure} and the interface-length optimizer: computation time and $\dot{H}^{-1}$ mix-norm decay rate.}
\label{tab:optimizer_comparison}
\end{table}

The results show that the interface-length optimizer requires less computation time while achieving comparable or higher $\dot{H}^{-1}$ mix-norm decay rates than the Sobolev-norm optimizer in all four settings. For the cellular-flow experiments, the interface-length method reduces the computation time by a factor of approximately $13$--$21$; for the Doswell experiments, the reduction factor is approximately $2.7$--$2.9$. In all four tested settings, when the two optimizers are run under the matched configuration described above and their respective controls are subsequently evaluated on the same reference transport solver, the control designed on the reduced model produces faster mix-norm decay: the measured $\dot{H}^{-1}$ decay rate improves by $3\%$--$64\%$ depending on the setting.

The speedup arises primarily from the lower per-iteration cost. Each iteration of the interface-length optimizer consists of a forward ODE solve ($O(N_p M)$ operations), an adjoint ODE solve of the same cost, a gradient evaluation, and a line search. By comparison, each iteration of the Sobolev-norm optimizer requires a forward PDE solve ($O(N_{\mathrm{cell}}^2 M)$), an adjoint PDE solve of equal cost, an elliptic solve for the $\dot{H}^{-1}$ norm, and the same gradient and line-search steps. The dominant cost difference lies in the forward and adjoint solves: tracking $N_p = 10^5$ markers along ODEs is cheaper than evolving a scalar field on a $500 \times 500$ grid through a full transport PDE. In our experiments, the interface model produces controls of comparable or better quality than the PDE-based optimizer at a fraction of the cost.


\subsection{Influence of the number of basis functions}\label{sec:N4}

The preceding experiments all use $N=2$ cellular-flow modes. A natural question is whether enlarging the control space to $N=4$ modes, $h_k(x_1,x_2)=\sin(k\pi x_1)\sin(k\pi x_2)$ for $k\in\{1,2,3,4\}$, yields improved mixing. We perform the optimization on $\Omega=(0,1)^2$ with $t\in[0,1]$, $M=1000$ time steps, $N_p=500{,}000$ markers, $\gamma_k=10^{-5}$ for all $k$, and constant initial guess $u_1=u_2=1$, $u_3=u_4=0$. The terminal ratio is $\|u^*\|_{L^2}/R_u = 0.15$ (projection inactive).

Figure~\ref{fig:interface_Cellular-N4} shows the interface evolution. The flow generates strong stretching and filamentation, with patterns visibly more complex than the $N=2$ optimum (Figure~\ref{fig:interface_Cellular-1}). The Eulerian scalar field in Figure~\ref{fig:theta_Cellular-N4} exhibits a four-cell vortex structure. The quantitative results (Figure~\ref{fig:Cellular_N4_mixnorm-length}) decompose into three observations. First, the optimizer is successful at the level of its own criterion: the enlarged control space yields a strictly deeper minimum of the discrete cost $J_h$ than the $N=2$ run, and the interface length at $T=1$ reaches approximately $23{,}000$, about $16$ times larger than the $N=2$ value of $1{,}450$, consistent with the exponential upper bound \eqref{eq:length_u_bound}. Second, this geometric gain does not transfer proportionally to the true mixing objective: the $\dot{H}^{-1}$ mix-norm decay rate measured on the reference transport solve is only $1.58$, compared to $2.59$ for $N=2$. Third, the gap between these two observations is the methodological content of the experiment: by Proposition~\ref{prop:zero_contour} and Remark~\ref{rem:critical_mixnorm}, $\mathcal L$ is a Lagrangian/BV-type indicator on $\Gamma_t$ while $\dot{H}^{-1}$ is a global Eulerian one, and the two can genuinely separate when the optimizer drives the interface into fine-scale structures localized within coherent regions.

\begin{figure}[ht]
  \centering
  \begin{subfigure}[b]{0.16\textwidth}
    \includegraphics[width=\textwidth]{./images/Cellular-N4/interface_t0.jpg}
    \caption{$t=0$}
  \end{subfigure}
  \begin{subfigure}[b]{0.16\textwidth}
    \includegraphics[width=\textwidth]{./images/Cellular-N4/interface_t1.jpg}
    \caption{$t=0.2$}
  \end{subfigure}
  \begin{subfigure}[b]{0.16\textwidth}
    \includegraphics[width=\textwidth]{./images/Cellular-N4/interface_t2.jpg}
    \caption{$t=0.4$}
  \end{subfigure}
  \begin{subfigure}[b]{0.16\textwidth}
    \includegraphics[width=\textwidth]{./images/Cellular-N4/interface_t3.jpg}
    \caption{$t=0.6$}
  \end{subfigure}
  \begin{subfigure}[b]{0.16\textwidth}
    \includegraphics[width=\textwidth]{./images/Cellular-N4/interface_t4.jpg}
    \caption{$t=0.8$}
  \end{subfigure}
  \begin{subfigure}[b]{0.16\textwidth}
    \includegraphics[width=\textwidth]{./images/Cellular-N4/interface_t5.jpg}
    \caption{$t=1.0$}
  \end{subfigure}
  \caption{Interface evolution under the optimized $N=4$ cellular flow.}
  \label{fig:interface_Cellular-N4}
\end{figure}

\begin{figure}[ht]
  \centering
  \begin{subfigure}[b]{0.16\textwidth}
    \includegraphics[width=\textwidth]{./images/Cellular-N4/theta_t0.jpg}
    \caption{$t=0$}
  \end{subfigure}
  \begin{subfigure}[b]{0.16\textwidth}
    \includegraphics[width=\textwidth]{./images/Cellular-N4/theta_t1.jpg}
    \caption{$t=0.2$}
  \end{subfigure}
  \begin{subfigure}[b]{0.16\textwidth}
    \includegraphics[width=\textwidth]{./images/Cellular-N4/theta_t2.jpg}
    \caption{$t=0.4$}
  \end{subfigure}
  \begin{subfigure}[b]{0.16\textwidth}
    \includegraphics[width=\textwidth]{./images/Cellular-N4/theta_t3.jpg}
    \caption{$t=0.6$}
  \end{subfigure}
  \begin{subfigure}[b]{0.16\textwidth}
    \includegraphics[width=\textwidth]{./images/Cellular-N4/theta_t4.jpg}
    \caption{$t=0.8$}
  \end{subfigure}
  \begin{subfigure}[b]{0.16\textwidth}
    \includegraphics[width=\textwidth]{./images/Cellular-N4/theta_t5.jpg}
    \caption{$t=1.0$}
  \end{subfigure}
  \caption{Scalar field $\theta_h$ under the optimized $N=4$ cellular flow.}
  \label{fig:theta_Cellular-N4}
\end{figure}

\begin{figure}[H]
  \centering
  \begin{subfigure}[b]{0.32\textwidth}
    \centering
    \includegraphics[width=\textwidth]{./images/Cellular-N4/interface_length.png}
    \caption{Interface length $\mathcal{L}(t)$.}
    \label{fig:Cellular_N4_length}
  \end{subfigure}
  \begin{subfigure}[b]{0.32\textwidth}
    \centering
    \includegraphics[width=\textwidth]{./images/Cellular-N4/mixnorm.png}
    \caption{$\dot{H}^{-1}$ mix-norm.}
    \label{fig:Cellular_N4_mixnorm}
  \end{subfigure}
  \begin{subfigure}[b]{0.32\textwidth}
    \centering
    \includegraphics[width=\textwidth]{./images/Cellular-N4/u1u2u3u4.png}
    \caption{Optimal $u_1,\dots,u_4$.}
    \label{fig:Cellular_N4_u1u2u3u4}
  \end{subfigure}
    \caption{Quantitative results under the optimized $N=4$ cellular flow.}
  \label{fig:Cellular_N4_mixnorm-length}
\end{figure}

The controls (Figure~\ref{fig:Cellular_N4_u1u2u3u4}) reveal the mechanism behind this behavior. The higher-frequency modes $u_3$ and $u_4$ are active only during the initial transient $t< 0.5$, after which they decay to zero and the controls revert to the low-mode $N=2$ behavior. The optimizer thus discovers that the additional modes are beneficial only for short-time stirring; sustained use of the higher harmonics creates fine-scale vortex cells that trap the scalar locally rather than transporting it across the domain.

Concretely, the local-vs-global mismatch noted in Remark~\ref{rem:local_global} explains the gap: the higher-frequency modes $u_3, u_4$ concentrate the additional stretching inside small-scale vortex cells where the scalar is trapped locally, so a much longer interface does not translate into faster global $\dot{H}^{-1}$ decay.


\section{Conclusions and perspectives}\label{sec:conclusions}

\subsection{Conclusions}
We proposed a reduced-order optimal control framework whose objective is the length of an advected material interface. Within a finite-dimensional smooth stream-function ansatz we established existence of a minimizer (Theorem~\ref{prop:H3}), derived the continuous adjoint and a discrete-consistent gradient (Remark~\ref{rem:discrete_exact}), and showed on cellular and Doswell benchmarks that the optimizer is faster than the Eulerian Sobolev-norm optimizer of \cite{hu2026structure} while producing exponential-type interface and mix-norm decay. The $N=4$ experiment shows the surrogate is useful but not fully faithful, in line with Glowinski's two-grid philosophy~\cite{Glowinski1992} that optimization on a well-chosen reduced model can be a useful companion to direct PDE-based optimization.

\subsection{Perspectives}
Several directions for future work arise.
\begin{itemize}
\item \textbf{Beyond interface length: sharper Lagrangian surrogates.}
A natural direction is to characterize the $\dot{H}^{-1}$-mixing functional more sharply in terms of the evolving interface, complementing partial results on exponential mixing rates \cite{alberti2016exponential, crippa2017cellular, brue2024enhanced}, and to design adaptive control-dimension strategies motivated by the non-monotone $N=4$ behavior. One concrete refinement replaces the local length functional by a non-local interface-supported $\dot{H}^{-1}$ energy of the line measure $\mu_{\Gamma_t}=\mathcal H^1\!\lfloor_{\Gamma_t}$ (Remark~\ref{rem:critical_mixnorm}), rewarding both elongation and spatial dispersion.

\item \textbf{Mixed fine-forward/reduced-adjoint strategy.} In the spirit of Glowinski's two-grid philosophy~\cite{Glowinski1992}, evolving the scalar on a fine Eulerian grid while computing adjoints on the reduced Hamiltonian model would couple PDE accuracy to the cheap reduced adjoint.
\end{itemize}




\section*{Acknowledgments}
The authors thank Yongcun Song for helpful discussions.


\bibliographystyle{abbrv}
\bibliography{refs_mixing.bib}

\end{document}
